\chapter{分词——教机器「认字」}
\label{ch:tokenization}
%% ============================================================

在模型能够处理文本之前，
必须将自然语言文本转换成模型可以操作的数字序列。
这个过程叫做\textbf{分词}（Tokenization）。
分词看似只是一个预处理步骤，
但它对模型的性能、效率和多语言能力有着深远的影响：
一个糟糕的分词器可以让一个优秀的模型变得平庸。

要理解分词的重要性，考虑一个简单的类比。
人类阅读时，不是逐字母扫描的：
我们以「词」为单位进行视觉识别，
熟练读者甚至可以一眼扫过整个短语。
如果强迫你逐字母阅读（就像字符级分词），
你的阅读速度和理解能力都会大幅下降。
反过来，如果给你一本用摩尔斯电码写的书
（每个词都是一个独立的长编码），
你同样会读不下去，因为「词汇」太多，记不住。
分词的任务就是在这两个极端之间找到最优的阅读单元：
既不太细（导致效率低下），
也不太粗（导致无法处理未知内容）。

对于 LLM 来说，分词的影响更加直接和深刻：
\textbf{分词决定了模型「看到」的输入长什么样}。
同一段文本，不同的分词方式会产生完全不同的 token 序列：
模型在此基础上学到的表示也会不同。
一个对中文不友好的分词器，
会让中文文本消耗更多的 token，
从而压缩模型的有效上下文窗口：
这等价于给模型戴上了「中文近视眼镜」。

\section{从词向量到子词：一段简史}

在讨论现代分词方法之前，
有必要回顾一下 NLP 领域是如何一步步走到今天的。
这段历史揭示了一个核心矛盾：
如何在词汇覆盖率和语义表达力之间找到平衡。

\subsection{词向量的黄金时代}

2013 年，Mikolov 等人提出了 Word2Vec，
第一次大规模地证明了
\textbf{词的语义可以被编码在低维向量中}。
Word2Vec 的核心思想简洁而深刻：
「你可以通过一个词的邻居来了解它。」
，一个词的含义由它频繁共现的词决定（分布假设）。
% NEW REF: mikolov2013word2vec = Mikolov, Tomas and Chen, Kai and Corrado, Greg and Dean, Jeffrey. Efficient Estimation of Word Representations in Vector Space. arXiv preprint arXiv:1301.3781, 2013.

Word2Vec 训练出来的向量展现了令人惊叹的代数性质：
$\vec{\text{King}} - \vec{\text{Man}} + \vec{\text{Woman}} \approx \vec{\text{Queen}}$。
这个结果让整个 NLP 社区为之兴奋：
词向量似乎捕获了某种深层的语义结构。

2014 年，Stanford 的 Pennington 等人提出了 GloVe
（Global Vectors for Word Representation），
通过对全局词共现矩阵做分解来学习词向量。
GloVe 的优势在于它同时利用了局部上下文
（像 Word2Vec）和全局统计信息（像传统的矩阵分解方法）。
% NEW REF: pennington2014glove = Pennington, Jeffrey and Socher, Richard and Manning, Christopher D. GloVe: Global Vectors for Word Representation. In Proceedings of EMNLP, 2014.

然而，这些\textbf{静态词向量}有一个根本性的缺陷：
每个词只有一个固定的向量表示，
不管它出现在什么上下文中。
“bank” 在「river bank」和「investment bank」中
得到完全相同的向量：
这显然丢失了关键的语义区分信息。

\subsection{上下文嵌入的崛起}

2018 年，Allen AI 的 Peters 等人提出了 ELMo
（Embeddings from Language Models），
用双向 LSTM 语言模型为每个词生成
\textbf{依赖上下文的}向量表示。
同一个 “bank” 在不同句子中会得到不同的嵌入。
这是一个范式性的转变：
从静态的字典式表示，走向动态的上下文式表示。
% NEW REF: peters2018elmo = Peters, Matthew E and Neumann, Mark and Iyyer, Mohit and Gardner, Matt and Clark, Christopher and Lee, Kenton and Zettlemoyer, Luke. Deep contextualized word representations. In Proceedings of NAACL-HLT, 2018.

随后，BERT~\cite{devlin2019bert} 将上下文嵌入推向了新的高度：
用 Transformer 替代 LSTM，在大规模语料上预训练，
然后微调到下游任务。
BERT 的成功不仅仅在于更好的嵌入质量，
还在于它证明了一个重要的观点：
\textbf{分词和嵌入可以被解耦}。
模型不需要以完整的词作为输入单元，
子词（subword）同样可以工作，甚至工作得更好。

BERT 使用的 WordPiece 分词器首次在大规模预训练中证明了子词分词的可行性。
它的词表只有 30,522 个 token，远小于词级别词表的数十万条目，
但在 11 项 NLP 基准上取得了当时的最佳成绩。
这个结果传达了一个强有力的信息：
\textbf{表示的粒度不需要与语言学的「词」对齐}。
模型可以从子词级别的输入中学到等价甚至更好的语义表示，
同时享受更小词表带来的效率优势。

这一洞察对后来的 GPT 系列至关重要。
GPT-2 和 GPT-3 都使用了 BPE 分词器：
正是因为 BERT 已经证明了子词分词在大模型上的有效性，
OpenAI 才有信心将 BPE 作为 GPT 系列的标准分词方案。
可以说，\textbf{现代 LLM 的分词选择在 BERT 时代就已经被奠定了}。

\subsection{未登录词问题}

从 Word2Vec 到 ELMo，所有基于词级别的方法
都面临一个致命的问题：\textbf{未登录词}（Out-of-Vocabulary, OOV）。

词表是在训练时固定的。
如果测试时遇到了训练语料中没有出现过的词
（例如新出现的专有名词 “ChatGPT”、
技术术语 “quantization”、
或者拼写错误 “langauge”），
模型只能将其映射为一个通用的 \texttt{<UNK>} token，
丢失所有语义信息。

英语的形态变化加剧了这个问题：
“play”、“plays”、“played”、“playing”、“player”、“playful”
，这些词共享词根 “play”，语义高度相关，
但在词级别的词表中它们是完全独立的条目。
词表要么包含所有变体（导致词表巨大），
要么遗漏某些变体（导致 OOV）。

对于中文、日文、韩文等语言，
问题又有不同的表现形式。
中文没有明确的词边界：
「研究生命的起源」可以被切分为
「研究/生命/的/起源」或「研究生/命/的/起源」：
分词歧义本身就是一个经典的 NLP 难题。

韩语虽然使用了空格，但一个韩语「词」（어절）通常由词根加多个助词组成（类似日语的活用形式），一个韩语「词」在语义上可能对应英语的一个完整短语。德语则以其复合词著称，如“Donaudampfschifffahrtsgesellschaftskapitän”（多瑙河蒸汽船运输公司船长）在德语中是一个合法的词，在词级别词表中需要作为独立条目收录。

这些跨语言的多样性说明了一个根本问题：\textbf{不存在一种适用于所有语言的「词」的定义}。子词分词的魅力正在于此，它绕过了「什么是词」这个语言学上的难题，转而从数据中自动发现合适的切分单元。

\textbf{子词分词}（subword tokenization）正是为解决这些问题而诞生的。
它的核心洞察是：
不要试图穷举所有的词，
而是找到一组合适的\textbf{子词单元}，
使得任何词都可以由这些子词拼接而成。

\section{为什么不能直接用字或词？}

让我们更系统地分析三种切分粒度的优劣。

\textbf{字符级切分}：
每个字母（或汉字）作为一个 token。
词表极小（英语仅需 $\sim$256 个字节），
永远不会遇到 OOV 问题。
但代价是巨大的：
字符层面的信息太稀疏了：
模型需要从大量的字符序列中学习词的含义，
效率极低。而且对于长文本，字符序列会非常长，
大幅增加注意力机制的计算成本。
以英语为例，一段 1000 词的文本
在字符级别约有 5000 个 token，
而在子词级别只有约 1300 个。
注意力的 $O(n^2)$ 复杂度意味着字符级的计算量是子词级的约 15 倍。

\textbf{词级切分}：
「I love coding」$\to$ [“I”, “love”, “coding”]。
问题更严重：英语有数十万个词，
而且新词不断出现（如 “ChatGPT”、“DeepSeek”）。
按词切分的词表要么巨大无比，要么遇到未登录词就束手无策。
对于中文、日文等语言，「什么是一个词」本身就是模糊的。
更实际的问题是：
巨大的词表意味着巨大的嵌入矩阵：
一个 50 万词的词表配上 4096 维嵌入就需要 2B 参数，
这比许多小模型的全部参数还多。

\textbf{子词级切分}，即 BPE 等方法所采用的：
在两个极端之间找到了最优的平衡：
高频词保留为完整的 token，
低频词或未见过的词被拆分为子词片段。
词表大小可控（32K--150K），
同时保证零 OOV。

\subsection{粒度选择的信息论视角}

从信息论的角度看，分词粒度的选择实质上是一个\textbf{编码效率}问题。Shannon 的信源编码定理告诉我们，最优编码应该使得每个符号携带的平均信息量接近信源的熵~\cite{shannon1948}。

字符级切分的问题在于：每个字符携带的信息量太少。英文字母 “t” 的信息量约为 $-\log_2(0.09) \approx 3.5$ bits，但一个英文子词（如 “the”）的信息量可能高达 8--12 bits。当模型以字符为单位处理时，它需要在注意力机制中跨越更长的距离来组合这些低信息密度的符号，这不仅增加了计算量，更增加了学习的难度。

词级切分则走向了另一个极端：每个词携带的信息量很高，但符号空间太大。一个 50 万词的词表意味着 softmax 层需要在 50 万个候选中做选择，这既增加了计算成本，也使得低频词的嵌入得不到充分训练。

子词切分的妙处在于：它通过数据驱动的方式自动找到信息密度与词表大小之间的最优平衡点。高频模式（如 “the”、“ing”、“tion”）被压缩为单个 token，每个 token 携带适中的信息量；低频模式则被分解为更细粒度的子词，确保覆盖率。BPE 的合并过程本质上就是在执行一种贪心的 Huffman 编码变体，优先为高频模式分配短码。

\subsection{子词分词的「形态学」效应}

子词分词的一个意外收获是：它在某种程度上恢复了语言的\textbf{形态学结构}，尽管它完全没有被显式地教导任何语言学知识。

以英语为例。BPE 训练在大规模英语语料上之后，通常会自动发现以下子词单元：
\begin{itemize}[noitemsep]
  \item 前缀：“un”、“re”、“pre”、“dis”
  \item 后缀：“ing”、“tion”、“ness”、“able”、“ment”
  \item 词根：“play”、“work”、“think”
\end{itemize}

这些子词单元恰好对应了英语的形态学组件。“unplayable” 可能被分为 \texttt{[un, play, able]}，完美地捕获了否定前缀、词根和形容词后缀。这意味着模型在处理从未见过的词时（如 “uncodable”），可以从已知的子词（“un”、“cod” 或 “code”、“able”）中推断其含义。

对于形态学更丰富的语言（如芬兰语、土耳其语），这个效应更加显著。芬兰语的一个词可以包含多达 7--8 个形态学后缀（如 “talonissakkaanko” = “在他们的房子里也是吗？”）。词级别的分词需要为这种高度组合的形态体系维护一个天文数字的词表，但 BPE 可以自然地将其分解为组成部分~\cite{sennrich2016bpe}。

这种「无监督形态学发现」是 BPE 成功的一个重要但常被忽视的因素。它解释了为什么基于子词的模型在形态学丰富的语言上大幅优于基于词的模型，不是因为模型更大或数据更多，而是因为分词器提供了更合理的输入表示。

\section{BPE：折中的智慧}

\subsection{从数据压缩到自然语言}

字节对编码（Byte Pair Encoding, BPE）的历史
比大多数人想象的要久远得多。
它最初由 Philip Gage 在 1994 年提出，
是一种简单而高效的\textbf{数据压缩}算法：
反复找到数据中出现最频繁的字节对，
用一个新的字节替换它们，从而缩短数据长度。
% NEW REF: gage1994bpe = Gage, Philip. A New Algorithm for Data Compression. C Users Journal, 12(2):23--38, 1994.

2016 年，Sennrich 等人~\cite{sennrich2016bpe}
将这个想法引入了自然语言处理，
用于解决机器翻译中的稀有词问题。
从此，BPE 成为了几乎所有现代 LLM 的标准分词方法。

这个类比值得深思：
BPE 本质上是在对语言做\textbf{压缩}：
将高频模式编码为单个 token，
低频模式则由多个 token 组合表示。
这与 Shannon 的信息论完美契合：
最优编码应该给高频符号分配短码，
给低频符号分配长码。
BPE 正是在做这件事：
只不过它的「码」是 token，
它的「符号」是自然语言中的子词片段。

\subsection{算法详解}

BPE 的算法流程如下：
\begin{enumerate}[noitemsep]
  \item 初始词表 = 所有单个字符（或字节）。
  \item 在训练语料中统计所有相邻 token 对的频率。
  \item 合并频率最高的一对，形成新的 token 加入词表。
  \item 重复步骤 2--3，直到词表大小达到目标（通常 32K--150K）。
\end{enumerate}

\subsection{完整的合并示例}

让我们用一个具体的例子来完整演示 BPE 的合并过程。
假设训练语料中频繁出现四个词：
“low”（5 次）、“lower”（2 次）、
“newest”（6 次）、“widest”（3 次）。

首先，每个词被拆成字符序列，词尾加上特殊标记 \texttt{\_}
（表示词边界）。初始表示如下：

\begin{center}\small
\renewcommand{\arraystretch}{1.3}
\begin{tabular}{ll}
\toprule
\rowcolor{figLightGray}
\textbf{词（频次）} & \textbf{Token 序列} \\
\midrule
low (5)    & \texttt{l o w \_} \\
lower (2)  & \texttt{l o w e r \_} \\
newest (6) & \texttt{n e w e s t \_} \\
widest (3) & \texttt{w i d e s t \_} \\
\bottomrule
\end{tabular}
\end{center}

\textbf{第 1 轮}：统计所有相邻 token 对的频率。
\texttt{(e, s)} 出现 $6 + 3 = 9$ 次（来自 newest 和 widest），
频率最高。合并 \texttt{e + s} $\to$ \texttt{es}。

\begin{center}\small
\renewcommand{\arraystretch}{1.3}
\begin{tabular}{ll}
\toprule
\rowcolor{figLightGray}
\textbf{词（频次）} & \textbf{Token 序列} \\
\midrule
low (5)    & \texttt{l o w \_} \\
lower (2)  & \texttt{l o w e r \_} \\
newest (6) & \texttt{n e w \underline{es} t \_} \\
widest (3) & \texttt{w i d \underline{es} t \_} \\
\bottomrule
\end{tabular}
\end{center}

\textbf{第 2 轮}：\texttt{(es, t)} 出现 $6 + 3 = 9$ 次，
频率最高。合并 \texttt{es + t} $\to$ \texttt{est}。

\begin{center}\small
\renewcommand{\arraystretch}{1.3}
\begin{tabular}{ll}
\toprule
\rowcolor{figLightGray}
\textbf{词（频次）} & \textbf{Token 序列} \\
\midrule
low (5)    & \texttt{l o w \_} \\
lower (2)  & \texttt{l o w e r \_} \\
newest (6) & \texttt{n e w \underline{est} \_} \\
widest (3) & \texttt{w i d \underline{est} \_} \\
\bottomrule
\end{tabular}
\end{center}

\textbf{第 3 轮}：\texttt{(est, \_)} 出现 $6 + 3 = 9$ 次，
频率最高。合并 \texttt{est + \_} $\to$ \texttt{est\_}。

\textbf{第 4 轮}：\texttt{(l, o)} 出现 $5 + 2 = 7$ 次，
频率最高。合并 \texttt{l + o} $\to$ \texttt{lo}。

\textbf{第 5 轮}：\texttt{(lo, w)} 出现 $5 + 2 = 7$ 次，
频率最高。合并 \texttt{lo + w} $\to$ \texttt{low}。

此时的状态：

\begin{center}\small
\renewcommand{\arraystretch}{1.3}
\begin{tabular}{ll}
\toprule
\rowcolor{figLightGray}
\textbf{词（频次）} & \textbf{Token 序列} \\
\midrule
low (5)    & \texttt{\underline{low} \_} \\
lower (2)  & \texttt{\underline{low} e r \_} \\
newest (6) & \texttt{n e w \underline{est\_}} \\
widest (3) & \texttt{w i d \underline{est\_}} \\
\bottomrule
\end{tabular}
\end{center}

继续合并下去，“low” 最终会被编码为单个 token \texttt{low\_}，
“newest” 会被编码为两个 token \texttt{new} + \texttt{est\_}，
以此类推。
这个过程完美地体现了 BPE 的核心思想：
\textbf{高频模式被压缩为单个 token，
低频模式由子 token 组合表示}。

\subsection{BPE 合并的可视化}

下面的图展示了 BPE 合并操作的核心机制：每一轮选出频率最高的相邻 token 对，将其合并为一个新 token，直到词表达到目标大小。

\begin{figure}[htbp]
\centering
\begin{tikzpicture}[
  tok/.style={
    rectangle, rounded corners=2pt,
    draw=figBlue!50, line width=0.6pt,
    fill=figBlue!8,
    minimum width=14mm, minimum height=9mm,
    inner sep=3pt, font=\small\ttfamily,
    text=figBlue!80!black
  },
  merged/.style={
    rectangle, rounded corners=2pt,
    draw=figOrange!60, line width=0.7pt,
    fill=figOrange!8,
    minimum width=14mm, minimum height=9mm,
    inner sep=3pt, font=\small\ttfamily\bfseries,
    text=figOrange!85!black
  },
  result/.style={
    rectangle, rounded corners=2pt,
    draw=figGreen!60, line width=0.7pt,
    fill=figGreen!10,
    minimum width=16mm, minimum height=9mm,
    inner sep=3pt, font=\small\ttfamily\bfseries,
    text=figGreen!85!black
  },
  roundbox/.style={
    rectangle, rounded corners=2pt,
    draw=figGray!30, fill=figGray!6,
    minimum width=14mm, minimum height=7mm,
    inner sep=3pt, font=\scriptsize\bfseries,
    text=figGray
  },
  arr/.style={
    -{Stealth[length=2.5pt,width=2.5pt]}, line width=0.6pt, figGray!60 % standardized
  },
  brace/.style={
    decorate, decoration={brace, amplitude=3pt, mirror},
    line width=0.6pt, figOrange!70
  },
]

% --- Round 1 ---
\node[roundbox] (r1label) at (0, 0) {Round 1};
\node[tok, right=10mm of r1label] (r1a) {n};
\node[tok, right=3mm of r1a] (r1b) {e};
\node[tok, right=3mm of r1b] (r1c) {w};
\node[merged, right=3mm of r1c] (r1d) {e};
\node[merged, right=3mm of r1d] (r1e) {s};
\node[tok, right=3mm of r1e] (r1f) {t};
\node[tok, right=3mm of r1f] (r1g) {\_};

% Brace below merged pair
\draw[brace] (r1d.south west) -- (r1e.south east);
\node[font=\tiny, text=figOrange!80!black] at ($(r1d.south)!0.5!(r1e.south) + (0, -5mm)$) {merge \texttt{e\,+\,s}};

% --- Round 2 ---
\node[roundbox] (r2label) at (0, -20mm) {Round 2};
\node[tok, right=10mm of r2label] (r2a) {n};
\node[tok, right=3mm of r2a] (r2b) {e};
\node[tok, right=3mm of r2b] (r2c) {w};
\node[merged, right=3mm of r2c] (r2d) {es};
\node[merged, right=3mm of r2d] (r2e) {t};
\node[tok, right=3mm of r2e] (r2f) {\_};

% Brace below merged pair
\draw[brace] (r2d.south west) -- (r2e.south east);
\node[font=\tiny, text=figOrange!80!black] at ($(r2d.south)!0.5!(r2e.south) + (0, -5mm)$) {merge \texttt{es\,+\,t}};

% --- Round 3 ---
\node[roundbox] (r3label) at (0, -40mm) {Round 3};
\node[tok, right=10mm of r3label] (r3a) {n};
\node[tok, right=3mm of r3a] (r3b) {e};
\node[tok, right=3mm of r3b] (r3c) {w};
\node[result, right=3mm of r3c] (r3d) {est\_};

% Final result annotation
\node[font=\scriptsize, text=figGreen!80!black, right=8mm of r3d, align=left]
  {\textrm{Final:} \texttt{[n, e, w, est\_]}};

% --- Vertical progression arrows (pointing to merged result in next round) ---
\draw[arr] ($(r1d.south)!0.5!(r1e.south) + (0, -8mm)$) -- (r2d.north);
\draw[arr] ($(r2d.south)!0.5!(r2e.south) + (0, -8mm)$) -- (r3d.north);

\end{tikzpicture}
\caption{BPE 对”newest”的逐轮合并过程。每轮选取频率最高的相邻 token 对进行合并：
\texttt{e\,+\,s} $\to$ \texttt{es}（第~1 轮），
\texttt{es\,+\,t} $\to$ \texttt{est}（第~2 轮），
\texttt{est\,+\,\_} $\to$ \texttt{est\_}（第~3 轮）。
橙色 token 标记待合并对，绿色 token 表示最终压缩结果。}
\label{fig:bpe-merge}
\end{figure}

\subsection{分词时的贪心匹配}

训练阶段产出的是一个\textbf{有序的合并规则表}（merge table），记录了合并发生的先后顺序。分词（编码）阶段则使用这个合并规则表来切分新文本。

关键细节是：BPE 分词时按照合并规则的顺序，贪心地应用每条规则。以 “lowest” 为例（假设它不在训练语料中）：

\begin{enumerate}[noitemsep]
  \item 初始：\texttt{l o w e s t \_}
  \item 应用规则 1（\texttt{e+s} $\to$ \texttt{es}）：\texttt{l o w es t \_}
  \item 应用规则 2（\texttt{es+t} $\to$ \texttt{est}）：\texttt{l o w est \_}
  \item 应用规则 3（\texttt{est+\_} $\to$ \texttt{est\_}）：\texttt{l o w est\_}
  \item 应用规则 4（\texttt{l+o} $\to$ \texttt{lo}）：\texttt{lo w est\_}
  \item 应用规则 5（\texttt{lo+w} $\to$ \texttt{low}）：\texttt{low est\_}
\end{enumerate}

最终结果：\texttt{[low, est\_]}。模型从未见过 “lowest” 这个词，但 BPE 成功地将它切分为两个有意义的子词，即”low” 和 “est\_”（表示最高级后缀）。这正是 BPE 处理 OOV 问题的优雅之处：通过组合已知的子词片段来表示未知的词。

这个过程也揭示了 BPE 分词的一个特点：\textbf{分词结果依赖于合并规则的顺序}。如果合并顺序不同（例如先合并 \texttt{l+o} 再合并 \texttt{e+s}），中间状态会不同，但最终结果通常是相同的，因为 BPE 保证所有可以合并的 pair 最终都会被合并，只是顺序不同。

\subsection{代码实现}

\begin{lstlisting}[caption={BPE core algorithm (simplified)}]
def learn_bpe(corpus, vocab_size):
    # Start with character-level tokens
    vocab = set(all_characters_in(corpus))
    tokens = [list(word) for word in corpus]

    while len(vocab) < vocab_size:
        # Count all adjacent token pairs
        pair_counts = count_pairs(tokens)
        # Find the most frequent pair
        best_pair = max(pair_counts, key=pair_counts.get)
        # Merge that pair everywhere
        new_token = best_pair[0] + best_pair[1]
        vocab.add(new_token)
        tokens = merge_pair(tokens, best_pair, new_token)

    return vocab
\end{lstlisting}

这段代码展示了 BPE 的核心循环：
统计频率、找到最高频对、合并、重复。
\texttt{count\_pairs} 遍历所有 token 序列，
统计相邻 pair 的出现次数。
\texttt{merge\_pair} 将语料中所有该 pair 的出现替换为合并后的新 token。

\subsection{时间复杂度分析}

BPE 训练的时间复杂度值得分析。
假设语料有 $N$ 个 token，目标词表大小为 $V$。
每一轮合并需要：
(1) 遍历所有 token 对并统计频率，$O(N)$；
(2) 找到最高频对，$O(|\text{pairs}|)$，可用优先队列优化到 $O(\log|\text{pairs}|)$；
(3) 在语料中执行替换，最坏 $O(N)$。
总共需要 $V$ 轮合并，因此总复杂度为 $O(VN)$。

对于典型的规模（$V = 100\text{K}$，$N = 10^{10}$），
直接实现会非常慢。
实际的工业级实现（如 HuggingFace 的 \texttt{tokenizers} 库）
使用了大量优化技巧：
基于前缀树的高效合并、增量更新频率计数、
Rust 实现的底层运算。
这些优化使得在数十 GB 的语料上
训练一个 BPE 分词器只需要几个小时。

\subsection{BPE 的局限与反直觉行为}

BPE 虽然在实践中非常成功，但它并非完美。理解其局限性对于训练高质量的分词器至关重要。

\textbf{贪心合并的次优性}。BPE 的每一步合并都选择当前频率最高的 pair——这是一个贪心策略，不保证全局最优。考虑一个极端的例子：如果训练语料中 “abcabc” 频繁出现，BPE 可能先合并 \texttt{a+b} $\to$ \texttt{ab}，然后 \texttt{ab+c} $\to$ \texttt{abc}。但如果语料中同时有大量 “bcde”，那么先合并 \texttt{b+c} $\to$ \texttt{bc} 可能在全局上更优，因为 \texttt{bc} 既可以服务于 “abc” 也可以服务于 “bcde”。BPE 无法做出这种前瞻性的决策。

\textbf{分词边界的脆弱性}。BPE 的分词结果对输入文本的微小变化可能高度敏感。例如，“unhappy” 可能被分为 \texttt{[un, happy]}，但 “Unhappy”（首字母大写）可能被分为 \texttt{[U, nhappy]} 或 \texttt{[Un, happy]}，仅仅因为大小写的变化就产生了完全不同的分词结果。这种不一致性可能导致模型对大小写的处理不够鲁棒。

\textbf{频率偏见的固化}。BPE 的合并规则完全由训练语料的频率分布决定。如果训练语料以英语为主，那么英语的常见子词模式（如 “tion”、“ing”、“the”）会优先被合并为独立 token，而其他语言的常见模式可能在合并序列中排在很后面。这意味着\textbf{分词器的训练数据分布直接决定了哪些语言被「优待」}，一个在英语为主的语料上训练的 BPE 分词器，天然地对中文、阿拉伯语等语言不友好。

\textbf{token 边界与语义边界的不一致}。BPE 的合并完全基于频率，不考虑语义。一个常见的反直觉案例是：“tokenization” 可能被切分为 “token” + “ization”（语义合理）或 “tok” + “en” + “ization”（语义不合理），取决于合并顺序。虽然 Transformer 的多层注意力机制最终能从不同的子词组合中恢复词义，但不合理的切分增加了模型的学习负担，相当于给模型出了一道不必要的「拼图题」。

\textbf{上下文无关性}。BPE 是一个\textbf{上下文无关}（context-free）的算法——同一个字符串在任何上下文中都产生相同的分词结果。但自然语言中存在大量的上下文相关歧义：“lead” 在 “lead the team”（领导）和 “lead pipe”（铅管）中含义完全不同，理想情况下分词器应该给出不同的表示。当然，这种上下文相关的分词在 BPE 框架下无法实现，它需要模型本身在嵌入层之上解决歧义。这也是为什么 Transformer 的多层结构如此重要：底层负责将上下文无关的子词组合为有意义的词，高层负责基于上下文解决歧义。

\section{WordPiece 与 Unigram：其他选择}

BPE 并不是唯一的子词分词算法。
还有两种重要的替代方案：
Google 的 WordPiece 和 SentencePiece 中的 Unigram 模型。
三者的核心区别在于\textbf{合并策略}。

\subsection{WordPiece：基于似然的合并}

WordPiece 最初由 Google 为语音识别系统开发，
后来被 BERT~\cite{devlin2019bert} 采用而广为人知。
% NEW REF: schuster2012wordpiece = Schuster, Mike and Nakajima, Kaisuke. Japanese and Korean voice search. In ICASSP, 2012.

BPE 选择合并\textbf{频率最高}的 token 对，
而 WordPiece 选择合并后
\textbf{使语言模型似然提升最大}的 token 对。
具体来说，对于候选合并对 $(a, b) \to ab$，
WordPiece 计算的得分是：
\begin{equation}
  \text{score}(a, b) = \frac{\text{freq}(ab)}{\text{freq}(a) \times \text{freq}(b)}
\end{equation}
这个得分衡量的是 $a$ 和 $b$ 共同出现的频率
相对于它们各自独立出现的频率：
本质上是\textbf{互信息}（pointwise mutual information）的一种近似。

直觉上，BPE 倾向于合并\textbf{绝对频率}高的对，
而 WordPiece 倾向于合并\textbf{关联性}强的对。
例如，“t” + “h” 在英文中的绝对频率极高（因为 “th” 出现在大量词中），
BPE 会优先合并它。
但 “u” + “n” 虽然绝对频率不如 “th” 高，
它们共现的条件概率可能更高
（“un-” 作为前缀非常稳定），
WordPiece 可能会优先合并它。

\subsection{Unigram 模型：从大到小的剪枝}

Unigram 模型（由 Kudo 于 2018 年提出）
采用了与 BPE/WordPiece 完全相反的策略：
\textbf{从一个很大的初始词表开始，逐步删除 token}。
% NEW REF: kudo2018sentencepiece = Kudo, Tomas and Richardson, John. SentencePiece: A simple and language independent subword tokenizer and detokenizer for Neural Text Processing. In Proceedings of EMNLP, 2018.

初始词表通常包含所有字符和最频繁的子字符串
（可能有数十万个条目）。
然后，Unigram 反复计算每个 token 被移除后
对整体语言模型似然的影响：
移除后似然下降最少的 token 被剪掉。
这个过程重复直到词表缩小到目标大小。

Unigram 的一个独特优势是它支持\textbf{多种分词方式}。
对于同一段文本，可能存在多种合法的分词方案，
Unigram 模型为每种方案分配一个概率。
这允许在训练时使用\textbf{子词正则化}
（subword regularization）：
随机采样不同的分词方案，
增加训练数据的多样性，
类似于数据增强（data augmentation）的效果。

\subsection{Unigram 的数学基础}

Unigram 模型的数学框架比 BPE 更加优雅。它假设每个 token 的出现是独立的（unigram 假设），一段文本 $\mathbf{x}$ 的某种分词方案 $\mathbf{s} = (s_1, s_2, \ldots, s_m)$ 的概率为：

\begin{equation}
  P(\mathbf{s}) = \prod_{i=1}^{m} p(s_i)
\end{equation}

其中 $p(s_i)$ 是 token $s_i$ 在词表中的概率（通过 EM 算法从语料中估计）。最优分词方案就是使 $P(\mathbf{s})$ 最大化的那个，这可以通过 Viterbi 算法在 $O(n \cdot V_{\max})$ 时间内高效求解，其中 $n$ 是字符序列长度，$V_{\max}$ 是词表中最长 token 的长度。

子词正则化的实现也很自然：不总是选择概率最大的分词方案，而是按概率采样。例如，“unbelievable” 可能有多种分词方式：
\begin{itemize}[noitemsep]
  \item \texttt{[un, believ, able]} —— 概率 0.45
  \item \texttt{[un, believe, able]} —— 概率 0.30
  \item \texttt{[unbeliev, able]} —— 概率 0.15
  \item \texttt{[un, bel, iev, able]} —— 概率 0.10
\end{itemize}

在训练时随机采样这些方案，模型被迫学会从不同的子词组合中提取相同的语义，这种正则化效果类似于 dropout，增强了模型的鲁棒性~\cite{kudo2018sentencepiece}。

一个值得注意的细节是：子词正则化的强度可以通过温度参数 $\alpha$ 来控制。当 $\alpha \to \infty$ 时，总是选择概率最大的分词方案（等价于确定性分词）；当 $\alpha \to 0$ 时，所有合法的分词方案被等概率采样（最大正则化）。实际使用中，$\alpha = 0.1$ 到 $0.3$ 之间的值通常能在正则化效果和训练稳定性之间取得良好的平衡。

值得一提的是，BPE 版本的子词正则化也是可能的。BPE-Dropout~\cite{provilkov2020bpe} 通过在 BPE 合并过程中随机跳过某些合并步骤来实现类似的效果。但由于 BPE-Dropout 在主流 LLM 训练中未被广泛采用（可能是因为额外的复杂性和不确定的收益），这里不做详细展开。

\subsection{三种方法的对比}

\begin{center}\small
\renewcommand{\arraystretch}{1.3}
\begin{tabular}{lccc}
\toprule
\rowcolor{figLightGray}
\textbf{特征} & \textbf{BPE} & \textbf{WordPiece} & \textbf{Unigram} \\
\midrule
方向           & 自底向上（合并） & 自底向上（合并） & 自顶向下（剪枝） \\
合并依据       & 频率最高       & 似然提升最大     & 似然损失最小     \\
分词确定性     & 确定性         & 确定性           & 概率性           \\
子词正则化     & 不支持         & 不支持           & 原生支持         \\
代表模型       & GPT 系列       & BERT             & T5               \\
               & LLaMA, Mistral &                  &                  \\
               & Qwen           &                  &                  \\
\bottomrule
\end{tabular}
\end{center}

在实践中，三种方法在最终效果上的差异通常不大：
词表大小和训练数据的选择对性能的影响远大于合并策略的选择。
这也是为什么不同实验室选择不同算法的原因：
差异不大，所以选择更方便实现的那个就好。

一个有趣的经验观察是：在相同词表大小下，Unigram 倾向于生成更短的 token（因为剪枝策略保留了更多的短 token），而 BPE 倾向于生成更长的 token（因为合并策略鼓励不断组合）。这意味着 Unigram 的平均序列长度可能稍长于 BPE，但对罕见词的处理更细粒度。对于多语言模型，这个差异有时是有意义的，更细的粒度意味着更好的跨语言共享（不同语言可以共享更多的基础子词单元）。

\subsection{为什么 BPE 在 LLM 中占主导？}

尽管 Unigram 在理论上更优雅（概率性框架、子词正则化），但 BPE 在大规模语言模型中占据了绝对的主导地位：GPT 系列、LLaMA、DeepSeek、Mistral 都使用 BPE。这一现象背后有几个实际原因。

\textbf{实现的简洁性}。BPE 的训练算法本质上是一个循环：统计频率、合并、更新。这个算法极其容易并行化，不同的 worker 可以独立统计不同语料分片的 pair 频率，然后汇总。Unigram 模型的训练则涉及 EM 算法的迭代和 Viterbi 解码，实现复杂度更高，并行化更困难。在处理 TB 级别的训练语料时，实现的简洁性直接转化为训练速度的优势。

\textbf{分词的确定性}。BPE 的分词结果是确定性的，同一段文本总是产生相同的 token 序列。Unigram 的最优分词是确定的（Viterbi 解码），但其子词正则化功能（随机采样分词方案）引入了随机性。在分布式训练中，确定性的分词确保了不同 GPU 上相同数据的一致性，这简化了工程实现。

\textbf{生态系统的锁定效应}。GPT-2 首先使用了 BPE，后续的开源社区围绕 BPE 构建了大量工具和经验。当 LLaMA 选择分词方案时，选择 BPE（或其 SentencePiece 变体）意味着可以复用大量已有的工具和代码。这种生态系统的锁定效应使得 BPE 的主导地位不断自我强化。

\section{词表大小：一个微妙的权衡}

词表大小是一个重要的超参数。
不同模型的选择差异巨大，
背后反映了不同团队的设计哲学。

\begin{center}\small
\renewcommand{\arraystretch}{1.3}
\begin{tabular}{lrll}
\toprule
\rowcolor{figLightGray}
\textbf{模型} & \textbf{词表大小} & \textbf{分词方式} & \textbf{说明} \\
\midrule
GPT-2 (2019)       & 50,257   & BPE          & 早期 BPE，纯英文优化 \\
LLaMA 1 (2023)     & 32,000   & SentencePiece & 偏小，中文效率低 \\
LLaMA 2 (2023)     & 32,000   & SentencePiece & 同 LLaMA 1 \\
Mistral 7B (2023)  & 32,000   & SentencePiece & 同 LLaMA 1 \\
Gemma 2 (2024)     & 256,000  & SentencePiece & Google，超大词表 \\
GPT-4 (2023)       & 100,256  & tiktoken      & 多语言 + 代码 \\
LLaMA 3 (2024)     & 128,256  & tiktoken 风格 & 大幅扩展，多语言 \\
DeepSeek-V3 (2024) & 128,000  & BBPE          & 强化多语言和代码 \\
Gemma 3 (2025)     & 262,144  & SentencePiece & 目前最大的公开词表 \\
Qwen3 (2025)       & 151,643  & tiktoken 风格 & 中文优化，CJK 效率高 \\
LLaMA 4 (2025)     & 202,048  & tiktoken 风格 & 200 语言，多模态 token \\
\bottomrule
\end{tabular}
\end{center}

\begin{whybox}[从 32K 到 262K：词表大小演进的逻辑]
这张表揭示了一个清晰的趋势：
词表大小从 2023 年的 32K 迅速扩展到 2025 年的 150K--262K。
背后的驱动力有三个：
(1) \textbf{多语言}：LLaMA 4 声称支持 200 种语言，
需要足够大的词表来为每种语言提供合理的 fertility；
(2) \textbf{多模态}：LLaMA 4 的词表中
包含了图像 token（\texttt{<|image|>}）和其他视觉 token，
词表 ID 从 200,000 开始分配给特殊 token；
(3) \textbf{代码}：Rust、Mojo 等新语言的关键词和运算符
需要更大的词表来避免被拆分为多个子词 token。
Google 的 Gemma 3 更是直接使用 262K 的超大词表：
证明对于足够大的模型，嵌入层的额外开销可以被忽略。
\end{whybox}

\textbf{词表太小}：每个文本需要更多的 token 表示
（因为更多词被拆成子词），
输入序列变长，注意力计算更慢，
模型需要从更长的上下文中学习。
一个量化这个问题的指标是\textbf{fertility}
（生育率），即平均每个词被切分成多少个 token。
fertility 越接近 1，分词效率越高。

\textbf{词表太大}：嵌入矩阵（$V \times d$）消耗大量参数和内存。
稀有 token 的嵌入得不到充分训练。
对于小模型，嵌入层可能占总参数的很大比例。

近年来的趋势是增大词表：
GPT-2 用 50K，LLaMA 1 用 32K，
LLaMA 3 增大到 128K，Qwen3 更是用到了 150K。
这是因为更大的词表能更高效地编码多语言文本
（每个中文词可以用更少的 token 表示），
而现代模型的参数量足够大，
嵌入层占比相对较小。

\subsection{词表大小对模型能力的非线性影响}

词表大小对模型能力的影响并非线性的，存在多个微妙的交互效应。

\textbf{推理能力与 token 粒度}。一个鲜少被讨论的现象是：分词粒度会影响模型的推理深度。考虑一个算术问题 “What is 347 + 689?”。如果 “347” 被编码为单个 token，模型只有一个 Transformer 层的前向传播来「处理」这个数字。但如果 “347” 被拆分为 “3”、“4”、“7” 三个 token，模型有三次前向传播的机会来理解这个数字的结构，每一位数字都经历了独立的注意力计算。这意味着\textbf{更细的分词粒度可能为模型提供更多的「计算步骤」来处理复杂信息}。这个观察部分解释了为什么 \texttt{split\_digits=True}（将每个数字字符独立为一个 token）能显著提升数学能力~\cite{mcleish2024abacus}。

\textbf{嵌入层的参数效率}。对于一个 7B 参数的模型，32K 词表对应的嵌入矩阵约为 $32\text{K} \times 4096 \approx 130\text{M}$ 参数（约占总参数的 1.9\%）；而 256K 词表对应约 1B 参数（约占 14\%）。对于小模型，这个比例不可忽视，词表过大意味着嵌入层「吃掉」了过多的参数预算，留给 Transformer 层的参数减少了。这就是为什么 LLaMA 1（7B）使用 32K 词表而不是更大的词表，在那个参数规模下，32K 是一个合理的权衡。但对于 70B 或更大的模型，嵌入层占比降至 1\% 以下，词表大小的约束大幅放松。

\textbf{softmax 瓶颈}。词表大小直接决定了输出层的 softmax 计算量。每个 token 位置的预测都需要在 $V$ 个候选中做 softmax，$V = 256\text{K}$ 意味着每个位置需要计算 256K 个 logit 值。虽然现代 GPU 的并行计算能力使这个计算不构成瓶颈，但它确实增加了内存带宽压力。一些研究~\cite{press2017tying} 通过权重共享（让输入嵌入矩阵和输出投影矩阵共享参数）来缓解这个问题，这也是 LLaMA 和 GPT-2 采用的策略。

\textbf{训练吞吐量的影响}。词表大小通过一个间接但重要的机制影响训练效率：更大的词表意味着更短的 token 序列（因为更多的内容可以用更少的 token 表示），而更短的序列意味着更高的训练吞吐量，因为注意力机制的 $O(n^2)$ 复杂度使得序列长度是计算量的主要决定因素。具体来说，如果将词表从 32K 增大到 128K 使得平均 token 序列长度缩短了 20\%，那么注意力计算量减少了约 36\%（$1 - 0.8^2 = 0.36$）。这种效率增益可能超过了嵌入层参数增加带来的额外开销，尤其对于大模型。

\textbf{词表大小与上下文窗口的交互}。词表大小和上下文窗口之间存在一种「等价关系」：增大词表可以在相同数量的 token 内容纳更多的原始文本，效果类似于增大上下文窗口。对于处理中文的模型尤其如此：将词表从 32K 增大到 150K 可能使中文的有效上下文窗口翻倍。从成本角度看，增大词表（一次性的训练成本增加）可能比增大上下文窗口（持续的推理成本增加）更经济。

\begin{whybox}[为什么大词表能帮助中文？]
在小词表（32K）下，
许多中文常用词会被拆成 2--3 个 token，
例如「人工智能」可能被切成「人工” + “智” + “能」。
而在 150K 词表下，
「人工智能」可以作为一个完整的 token 存在。
这不仅节省了序列长度（效率提升），
也让模型能直接学习这个词的完整语义
（而不是从子词中拼凑），
提升了模型对中文的理解能力。

具体的数字可以说明这个差异。
LLaMA 1（32K 词表）处理中文时，
平均 fertility 约为 2.5：
即每个中文词平均被切成 2.5 个 token。
同样的中文文本在 Qwen3（150K 词表）下，
fertility 降至约 1.2。
这意味着相同的上下文窗口可以容纳约 2 倍的中文内容：
对中文理解能力的提升是直接且显著的。
\end{whybox}

\begin{warnbox}[fertility 的隐含成本]
fertility 不仅影响理解能力，还直接影响\textbf{推理成本}。
API 服务按 token 计费：
如果你的语言 fertility 是 2.5（即同等内容需要 2.5 倍的 token），
那么你每次调用的成本也是 fertility 为 1 的语言的 2.5 倍。
对于使用 LLaMA 1 分词器的模型处理中文，
用户在成本上受到了系统性的「不公平对待」。
这也是为什么中文 AI 社区强烈推动大词表的原因之一。
\end{warnbox}

\section{多语言分词：一个被低估的难题}

词表大小只是多语言挑战的冰山一角。不同语言的文字系统在结构上有着根本性的差异，给统一的分词方案带来了深层挑战。

\subsection{CJK 生育率问题}

中日韩（CJK）语言面临一个独特的困境。以中文为例，每个汉字在 UTF-8 编码中占 3 个字节，而英文字母只占 1 个字节。这意味着在字节级别的 BPE 中，一个汉字的「基础成本」就是英文字母的 3 倍。

如果训练语料以英文为主（这在大多数早期模型中都是如此），BPE 的合并规则会优先为英文子词分配 token，因为英文子词的频率更高。中文汉字虽然单个频率不低，但组合（如「人工智能」）的频率被分散到了数以万计的双字词、四字词中。结果就是：

\begin{center}\small
\renewcommand{\arraystretch}{1.3}
\begin{tabular}{lll}
\toprule
\rowcolor{figLightGray}
\textbf{文本} & \textbf{LLaMA 1 (32K)} & \textbf{Qwen3 (151K)} \\
\midrule
“artificial intelligence” & 2 tokens & 2 tokens \\
「人工智能」                  & 6 tokens & 1 token \\
“machine learning”         & 2 tokens & 2 tokens \\
「机器学习」                  & 5 tokens & 1 token \\
“The quick brown fox”      & 4 tokens & 4 tokens \\
「今天天气不错」              & 9 tokens & 3 tokens \\
\bottomrule
\end{tabular}
\end{center}

这不仅是效率问题，更是能力问题。当「人工智能」被切成 6 个 token 时，模型需要在 6 个位置的注意力计算中「重新组装」这个词的语义。而当它是 1 个 token 时，模型可以直接在嵌入空间中学习其完整语义。对于复杂的中文短语（如「量子纠缠态的贝尔不等式违反」），token 数量的膨胀意味着模型的有效上下文窗口被压缩，同样的 4096 token 窗口，英文可以容纳约 3000 词，而中文可能只能容纳约 1200 词（使用 LLaMA 1 分词器时）。

\subsection{阿拉伯语与希伯来语的 RTL 挑战}

从右到左（Right-to-Left, RTL）书写的语言给分词带来了额外的复杂性。阿拉伯语的文字系统有几个独特特征：

\textbf{连写形态}。阿拉伯字母的形状取决于它在词中的位置——同一个字母有独立形、词首形、词中形和词尾形四种写法。BPE 在字节级别操作时，这些不同形式对应不同的 UTF-8 编码，可能被视为完全不同的「字符」，分词器无法识别它们其实是同一个字母。

\textbf{附加符号}（diacritics）。阿拉伯语的元音通常用附加在辅音字母上方或下方的小符号表示。这些符号在正式文本中通常省略（读者依靠上下文推断元音），但在宗教文本和教学材料中会保留。一个分词器需要能够处理有附加符号和无附加符号的两种形式，而 BPE 会将它们视为完全不同的字符序列。

\textbf{混合文本方向}。当阿拉伯语文本中嵌入英语单词或数字时，文本方向会在 RTL 和 LTR 之间切换。这种混合文本在技术文档中极为常见。分词器的预分词步骤如果不能正确处理方向切换，就会在错误的位置断词~\cite{rust2021good}。

这些 RTL 挑战并非不可解决，但它们需要分词器设计者对目标语言有深入的了解，或者至少有来自该语言社区的反馈。2024--2025 年，阿拉伯语 AI 社区推动了多个专门针对阿拉伯语优化的分词器项目，显著改善了阿拉伯语 LLM 的分词效率。类似的社区驱动努力也在印地语、泰语、韩语等语言中展开。

\subsection{泰语与越南语的声调标记挑战}

泰语和越南语为分词带来了另一类独特的挑战：\textbf{声调标记}。

越南语使用拉丁字母，但每个字母上可以有声调符号（如 ă、ã、ả、ạ、ắ）。一个越南语元音字母加上声调符号在 UTF-8 中可以有两种表示方式：NFC 形式（组合字符，1 个码点）和 NFD 形式（基字符 + 组合标记，2 个码点）。如果分词器的训练数据混合了两种 Unicode 规范化形式，同一个越南语词就会产生不同的分词结果，这不仅降低了效率，还可能导致模型对同一个词学到两套不同的表示。

泰语则是一种无空格语言（类似中文），但词的边界更加模糊，泰语的「词」概念本身在语言学上就有争议。更复杂的是，泰语的声调标记和元音标记可以出现在辅音的上方、下方、前方或后方，形成了复杂的二维布局，而 BPE 只能处理一维的序列。

这些挑战的共同教训是：\textbf{分词不是一个纯粹的技术问题，它深刻地依赖于对目标语言的文字系统和正字法规则的理解}。一个没有泰语专家参与设计的分词器，几乎不可能在泰语上达到好的效率。

\subsection{代码作为一种「语言」}

编程语言与自然语言有着根本不同的统计特性，这对分词提出了独特要求。

自然语言的核心特征是 Zipf 分布，少量高频词占据了大部分文本。代码则不同：\textbf{标识符}（变量名、函数名）具有长尾分布，大量标识符只出现一两次。更重要的是，这些标识符携带了关键的语义信息，\texttt{calculate\_total\_revenue} 和 \texttt{calc\_rev} 表达相同的含义，但对分词器来说它们是完全不同的 token 序列。

\textbf{缩进即语法}。Python 用空白缩进表示代码块层次。4 个空格的缩进和 8 个空格的缩进传达了不同的语义信息（不同的嵌套层次）。如果分词器将空格序列以任意方式切分（例如 8 个空格被切为 “3 个空格 + 5 个空格”），模型就难以从 token 序列中可靠地恢复缩进层次。

\textbf{运算符的密度}。代码中密集的运算符序列（如 C++ 的 “\texttt{>>}”、“\texttt{<<=}”、“\texttt{->*}”）在自然语言中极为罕见。以自然语言为主训练的分词器通常不会将这些运算符组合为独立 token，它们会被逐字符切分，增加序列长度并丢失语义完整性。

\textbf{字符串字面值}。代码中的字符串字面值本质上是「嵌入在代码中的自然语言」。一个理想的分词器应该对字符串内容和代码语法使用不同的分词策略，但 BPE 无法区分这两种上下文。

\subsection{日语的独特复杂性}

日语可能是分词面临的最复杂的自然语言。它同时使用四种文字系统：\textbf{平假名}（表音，用于语法词和送假名）、\textbf{片假名}（表音，用于外来词）、\textbf{汉字}（表意，从中文借入）和\textbf{罗马字}（拉丁字母，用于缩写和外来名称）。一个日语句子可能同时包含这四种文字系统，例如「AIモデルは人工知能の一種です」混合了罗马字（AI）、片假名（モデル）、汉字（人工知能、一種）和平假名（は、の、です）。

对于 BPE 分词器，这意味着同一个概念在日语中可能有多种表达形式，每种形式占用不同的 token。「コンピュータ」（片假名，外来词 computer）和「計算機」（汉字，传统日语词）表达相同的含义，但在 BPE 词表中是完全不同的 token 序列。这种一义多表的特点使得日语的有效词表需求远大于表面看到的，也解释了为什么 Gemma 3 选择了 262K 的超大词表。

\subsection{印地语和低资源语言的困境}

对于低资源语言（训练数据稀少的语言），分词问题更加严峻。以印地语（天城文书写）为例：印地语在互联网上的文本量不到英语的 1\%，但它是世界第三大使用语言（超过 6 亿使用者）。

在以英语为主训练的 BPE 分词器中，印地语的 fertility 通常高达 4--6，即一个印地语词平均需要 4--6 个 token。这意味着使用这种分词器的模型在处理印地语时，有效上下文窗口缩小为英语的 1/4 到 1/6。更严重的是，低 fertility 与模型在该语言上的低性能形成恶性循环：token 浪费导致该语言的训练效率低下，性能差又减少了用户使用该语言的动机，进而减少了训练数据的来源。

打破这个循环的一种方法是\textbf{语言适配}（language adaptation），在已有分词器的基础上，为目标语言添加高频子词 token。具体做法是：在目标语言的语料上训练一个独立的 BPE 分词器，然后将其中最高频的 $k$ 个 token 追加到已有词表中。这种方法不需要重新训练整个分词器，但需要对模型进行一轮额外的 continual pre-training 来学习新 token 的嵌入~\cite{cui2024llama}。

Chinese-LLaMA~\cite{cui2024llama} 是这种方法的典型案例。它在 LLaMA 1 的 32K 词表基础上追加了约 20K 的中文 token，形成了一个 49K 的混合词表。新增的中文 token 包括常用汉字、双字词和高频多字词（如「人工智能」「深度学习」）。这种扩展使得中文的 fertility 从原来的约 2.5 降至约 1.5，一个显著的改进，但仍不如从头训练的中文优化分词器（如 Qwen 的 1.2）。

这种方法的局限性在于：追加的 token 的嵌入质量取决于 continual pre-training 的数据量和质量。如果 continual pre-training 的数据量不够大，新 token 的嵌入可能不够稳定，这实际上就是 glitch token 问题的一种温和形式。经验法则是：每个新增 token 至少需要在 continual pre-training 中出现 $10^5$ 次以上才能形成稳定的嵌入。

\subsubsection{KazByte：哈萨克语的「tokenizer tax」与字节级适配}

低资源语言被多语言分词器征收的「税」
往往要等到一个具体语言的 fertility 被定量测量出来才会被严肃对待。
2026 年 3 月发表的 \textbf{KazByte}~\cite{tursynbek2026kazbyte}
是其中一个典型样本。
研究者发现 Qwen2.5-7B
把哈萨克语切成的 token 数
\textbf{远超英语等价文本}：
同样信息量的段落，
哈萨克语版本比英语版本多消耗数倍 token，
推理成本、上下文窗口利用率
都因此承受额外的「tokenizer tax」。

KazByte 的方案是
\textbf{冻结 Qwen 主体不动}，
在前面接一个小型 adapter：
adapter 把哈萨克语字节流
重新映射到 Qwen 的内部表示空间，
然后只对 attention 层做轻量微调。
整个流程不重训分词器、不扩词表，
也不需要大规模的哈萨克语预训练语料。
对哈萨克语用户来说，
这相当于在保留 Qwen 全部英语/中文能力的同时
取消了多语言分词器对哈萨克语的隐性歧视。

KazByte 的更广义启示是：
当多语言 tokenizer 对某个低资源语言天然不友好时，
最现实的修复路径不一定是「重训一个公平的分词器」
（这往往要求该语言社区拿出 frontier 级别的算力），
而是\textbf{在字节层接一个轻量适配器}，
让模型绕开 tokenizer 的偏见去重新读取该语言。
这与本章前文讨论的字节级建模思路在精神上完全一致：
字节是所有语言的公约数，
任何 tokenizer 偏见都可以在字节层被回滚。

\subsection{如何构建多语言友好的分词器}

面对上述挑战，现代多语言分词器采用了几种关键策略：

\textbf{均衡采样训练数据}。Qwen3 的分词器在训练时显式控制了不同语言的采样比例，确保中文、英文、代码各占足够的份额。这与模型预训练数据的语言比例不同，分词器的训练更强调覆盖率，而非反映真实分布。

\textbf{预留特殊 token 空间}。LLaMA 4 的 202K 词表中，前 200K 个 ID 分配给常规 token，最后 2048 个 ID 预留给特殊 token（包括多模态 token）。这种设计允许在不改变分词器的情况下添加新的特殊 token。

\textbf{语言特定的预分词规则}。tiktoken 风格的分词器使用正则表达式进行预分词，将文本先切分为「词」级别的片段，然后在每个片段内部做 BPE。不同的正则表达式可以适配不同语言的特点，例如，对中文文本可以按字符切分（每个汉字作为一个独立片段），而对英文文本则按空格和标点切分。

\textbf{跨语言共享子词}。多语言分词器设计的一个核心权衡是：不同语言之间应该共享多少子词？一方面，共享子词可以促进跨语言的知识迁移，例如，如果中文的「信息」和日语的「情報」共享相同的汉字子词，模型可能更容易学到它们之间的语义关联。另一方面，过度共享可能导致「语言干扰」，一个子词在不同语言中有不同的含义（如英语的 “die” 和德语的 “die”（定冠词）），共享同一个嵌入可能造成语义混淆。

Qwen3~\cite{qwen2025qwen3} 的分词器在这个问题上做了精心的设计：CJK 汉字在中文和日语之间共享（因为它们有相近的语义），但中文特有的简化字和日语特有的假名分别拥有独立的 token。这种「选择性共享」策略在保持跨语言迁移的同时，避免了不当的语义干扰。

\textbf{Unicode 规范化的统一}。在构建多语言分词器之前，必须对训练数据做一致的 Unicode 规范化。NFC（Canonical Decomposition, followed by Canonical Composition）是推荐的标准，它确保同一个字符总是用相同的码点表示。如果训练数据中混合了 NFC 和 NFD 形式，分词器可能将视觉上完全相同的文本编码为不同的 token 序列。这种不一致性在阿拉伯语、越南语和韩语中尤为常见，因为这些语言的文字系统大量使用组合字符（combining characters）。

\textbf{数字和标点的语言无关处理}。不同语言使用不同的数字系统（阿拉伯数字 0--9、东阿拉伯数字 ٠--٩、天城文数字 ०--९ 等）和标点符号（中文句号「。」vs 英文句号「.」、中文逗号「，」vs 英文逗号「,」）。一个多语言分词器需要决定是将不同系统的数字和标点统一为一组 token（更高效但可能丢失文化信息），还是分别编码（保留文化信息但增大词表）。大多数现代分词器选择了后者，保留每种数字系统和标点符号的独立 token，因为这些信息在翻译和文化适配任务中可能很重要。

\textbf{分词器的 A/B 测试}。在工业实践中，分词器的选择通常不是基于理论分析做出的，而是通过\textbf{大规模 A/B 测试}。具体做法是：用候选分词器分别训练小规模的代理模型（proxy model，通常 1B--3B 参数），在一组覆盖多语言、代码和推理的基准上评估，然后选择表现最优的分词器用于大规模训练。这种方法的成本远低于直接在 frontier 规模上尝试不同的分词器，一次代理模型的训练只需要几千美元，而一次 frontier 训练可能需要几百万美元。

\section{SentencePiece：工业级实现}

SentencePiece 是 Google 开源的分词工具~\cite{kudo2018sentencepiece}，
被 LLaMA、Qwen、DeepSeek、Mistral 等主流模型采用。
它不仅是一个分词算法的实现，
更是一个\textbf{完整的分词框架}：
同时支持 BPE 和 Unigram 两种算法，
提供训练、编码、解码的全套功能。

\subsection{语言无关性}

SentencePiece 的核心特点是将输入视为\textbf{原始 Unicode 字节流}，
不依赖任何预分词（pre-tokenization）步骤。

传统的 BPE 实现（如 GPT-2 的 tokenizer）
需要先按空格和标点将文本切分为「词」，
然后在词内部做 BPE 合并。
这种设计对英文友好（空格是天然的词边界），
但对中文、日文等没有空格的语言不够自然。

SentencePiece 直接在字节级别操作：
将所有文本编码为 UTF-8 字节序列，
然后在字节序列上学习 BPE（或 Unigram）合并规则。
这意味着它可以处理\textbf{任何语言}：
甚至是 emoji、数学符号和未知的 Unicode 字符：
因为一切最终都是字节。

\subsection{字节回退机制}

SentencePiece 的一个关键设计是\textbf{字节回退}
（byte fallback）机制。
当遇到词表中不存在的 Unicode 字符时，
SentencePiece 不会输出 \texttt{<UNK>}，
而是将该字符拆解为其 UTF-8 字节表示，
用特殊的字节 token（如 \texttt{<0xE4>}）来编码。

这意味着 SentencePiece 的分词结果
\textbf{永远不会包含未知 token}：
任何输入都能被编码和解码。
这对于处理用户输入至关重要，
因为用户可能输入任何字符：
生僻汉字、emoji、甚至是乱码。

\subsection{确定性与可复现性}

SentencePiece 的设计还带来一个重要的工程优势：
分词器的输出是\textbf{确定性的}：
同一段文本总是产生相同的 token 序列，
与操作系统、locale 设置和文本预处理无关。
这在分布式训练中至关重要，
因为不同 GPU 上的分词结果必须完全一致。

\subsection{训练与使用}

使用 SentencePiece 训练一个分词器非常简单：

\begin{lstlisting}[caption={Training a SentencePiece tokenizer}]
import sentencepiece as spm

# Train a BPE tokenizer with 32K vocab
spm.SentencePieceTrainer.train(
    input='training_corpus.txt',
    model_prefix='my_tokenizer',
    vocab_size=32000,
    model_type='bpe',          # or 'unigram'
    character_coverage=0.9995, # cover 99.95% of characters
    byte_fallback=True,        # enable byte fallback
    split_digits=True,         # split each digit separately
    num_threads=16,
)

# Use the trained tokenizer
sp = spm.SentencePieceProcessor(model_file='my_tokenizer.model')
tokens = sp.encode('Hello world!', out_type=str)
# ['_Hello', '_world', '!']
ids = sp.encode('Hello world!', out_type=int)
# [12345, 2048, 67]
text = sp.decode(ids)
# 'Hello world!'
\end{lstlisting}

注意 \texttt{character\_coverage=0.9995} 这个参数：
它指定词表应覆盖训练语料中 99.95\% 的字符。
剩下的 0.05\% 极稀有字符将通过字节回退处理。
对于以 CJK 字符为主的语料，
这个值通常需要设得更高（如 0.99995），
因为 CJK 字符集本身就很大。

\subsection{BBPE：字节级 BPE}

DeepSeek-V3~\cite{deepseekv3} 使用了 BBPE（Byte-level BPE），一种在 SentencePiece 基础上的变体。BBPE 的核心思想是：不从字符级别开始合并，而是从\textbf{字节级别}开始。

传统的 BPE（如 GPT-2）的初始词表是 Unicode 字符（数千个），而 BBPE 的初始词表只有 256 个字节。这带来了一个重要优势：初始词表的大小是固定的，不受语言种类的影响。无论训练语料包含多少种语言，BBPE 的起点都是相同的 256 个字节，然后通过合并操作，自动发现每种语言中的高频模式。

GPT-4 的 tiktoken 也采用了类似的字节级 BPE 思路，但它在字节级别之上添加了正则表达式预分词，先用正则表达式将文本切分为「词」级别的片段，然后在每个片段内部做字节级 BPE。这种两阶段设计结合了语言知识（正则表达式中的词边界规则）和数据驱动（BPE 的统计合并），在实践中表现优异。

\subsection{SentencePiece 的性能调优}

在工业实践中，SentencePiece 的训练有几个关键调优参数值得深入讨论。

\textbf{预分词正则表达式}。tiktoken 风格的分词器使用的预分词正则表达式对最终的分词质量有显著影响。GPT-4 使用的正则表达式如下（简化版）：

\begin{lstlisting}[caption={GPT-4 pre-tokenization regex (simplified)}, language={}]
'(?i:[sdmt]|ll|ve|re)    # English contractions
|[^\r\n\p{L}\p{N}]?\p{L}+  # Optional punct + letters
|\p{N}{1,3}               # 1-3 digit numbers
| ?[^\s\p{L}\p{N}]+[\r\n]* # Punctuation sequences
|\s*[\r\n]                 # Newlines
|\s+(?!\S)                 # Trailing whitespace
|\s+                       # Other whitespace
\end{lstlisting}

这个正则表达式精心处理了英语的缩写形式（“don't” $\to$ “don” + “'t”）、数字（每 1--3 位独立分组）和空白字符（换行符单独处理）。不同的正则表达式会产生不同的预分词结果，进而影响 BPE 的合并行为，这就是为什么不同模型使用相同的 BPE 算法却产生不同的词表。

\textbf{词表去重与规范化}。训练完成的词表中可能包含视觉上相同但 Unicode 编码不同的 token（如 “fi”(U+FB01, Latin Small Ligature FI) 和 “fi”(U+0066 U+0069)）。工业级实现通常会对词表做 Unicode NFC 规范化，合并这些重复项，否则模型可能对同一个视觉符号学到完全不同的嵌入。

\textbf{特殊字符的处理}。某些字符在分词器训练中需要特殊处理：零宽度空格（U+200B）、BOM（U+FEFF）、各种 Unicode 控制字符。如果这些字符出现在训练语料中但未被正确处理，它们可能被合并到正常 token 中，产生难以调试的 bug。最佳实践是在分词器训练前用显式的预处理步骤移除或规范化这些字符。

\section{特殊 Token 的设计}

现代 LLM 的词表不仅包含自然语言 token，
还包含一系列\textbf{特殊 token}，用于控制模型的行为。

\subsection{基础特殊 Token}

\begin{itemize}[noitemsep]
  \item \texttt{<bos>}（begin of sequence）：
        标记序列的开始。
  \item \texttt{<eos>}（end of sequence）：
        标记序列的结束，模型在生成时遇到它就停止。
  \item \texttt{<pad>}（padding）：
        用于将不同长度的序列填充到相同长度，
        以便组成 batch。
\end{itemize}

\subsection{对话模板 Token}

随着 LLM 从文本补全转向多轮对话，
对话模板（chat template）的设计变得至关重要。
不同模型使用不同的对话格式：

\textbf{ChatML 格式}（OpenAI/Qwen）：
\begin{lstlisting}[language={}]
<|im_start|>system
You are a helpful assistant.
<|im_end|>
<|im_start|>user
What is BPE?
<|im_end|>
<|im_start|>assistant
BPE stands for Byte Pair Encoding...
<|im_end|>
\end{lstlisting}

\textbf{LLaMA 3 格式}：
\begin{lstlisting}[language={}]
<|begin_of_text|>
<|start_header_id|>system<|end_header_id|>

You are a helpful assistant.
<|eot_id|>
<|start_header_id|>user<|end_header_id|>

What is BPE?
<|eot_id|>
<|start_header_id|>assistant<|end_header_id|>
\end{lstlisting}

这些特殊 token 在预训练阶段不存在：
它们是在 SFT（监督微调）阶段被添加到词表中的。
模型通过 SFT 学会这些 token 的含义：
\texttt{<|im\_start|>} 表示一条消息的开始，
\texttt{<|im\_end|>} 表示消息结束。

对话模板的设计看似是一个纯工程问题，但它对模型的行为有着微妙而深远的影响。

\textbf{模板一致性}。如果推理时使用的对话模板与 SFT 训练时使用的不一致（例如漏掉了某个特殊 token，或者使用了不同的角色标记），模型的行为可能显著退化，它可能忽略系统提示、无法正确区分用户和助手的角色、或者在多轮对话中丢失上下文。这是实际部署中最常见的「静默 bug」之一：一切看起来都在工作，但模型的输出质量不如预期。

\textbf{模板对推理能力的影响}。研究表明~\cite{zheng2024judging}，对话模板中特殊 token 的位置会影响模型的注意力分配。具体来说，紧跟在 \texttt{<|im\_start|>assistant} 之后的 token 会得到更高的注意力权重，模型学会了在这些位置「更加认真地生成」。这解释了为什么在某些模型中，在回答开头添加一个简短的「思考」前缀（如 “Let me think about this...”）可以提升回答质量，不是因为这句话本身有什么魔力，而是因为它为模型提供了额外的 token 位置来「热身」。

\textbf{不同模板之间的迁移}。当用户想将一个为 ChatML 格式训练的模型用于 LLaMA 格式的提示时，通常需要一轮额外的 SFT 来适配新模板。但如果直接「硬切」（不经过适配），模型的行为可能变得不可预测，因为新模板中的特殊 token 对模型来说是「未知的信号」。HuggingFace 的 \texttt{chat\_template} 系统试图通过统一的 Jinja2 模板语言来标准化这个过程，使得不同模型可以自动应用正确的对话格式。

\subsection{推理与工具调用 Token}

2024--2025 年的一个重要趋势是
为特定功能设计专用的特殊 token：

\begin{itemize}[noitemsep]
  \item \texttt{<think>} / \texttt{</think>}：
        DeepSeek-R1~\cite{deepseekr1} 和 Qwen3~\cite{qwen2025qwen3}
        用于标记推理过程的 token，
        推理内容包在这对标签内，不展示给用户。
  \item \texttt{<tool\_call>} / \texttt{</tool\_call>}：
        标记工具调用请求：
        模型输出结构化的函数调用信息。
  \item \texttt{<tool\_response>} / \texttt{</tool\_response>}：
        标记工具返回的结果，
        模型根据这些结果继续生成回复。
  \item \texttt{<|code|>}：
        某些模型用于标记代码块的开始，
        触发代码补全模式。
\end{itemize}

这些特殊 token 的嵌入在 SFT 阶段被学习，
使模型理解「到这里该停了」或
「接下来是推理过程」等控制信号。
它们的设计看似简单，
但直接影响了模型在多轮对话、工具调用、
推理模式等场景中的行为。

\subsection{特殊 Token 的嵌入初始化}

一个工程上的微妙问题是：新添加的特殊 token 的嵌入向量应该如何初始化？

最简单的方法是随机初始化，但这会导致新 token 的嵌入与已有 token 的嵌入分布不一致，可能引起训练不稳定。一种更好的做法是用已有 token 嵌入的\textbf{均值}来初始化新 token，这确保了新 token 的嵌入处于已有嵌入的「中心」，不会在 SFT 训练初期产生极端的梯度。

LLaMA 3 的做法更加精细：对于语义明确的特殊 token（如 \texttt{<|start\_header\_id|>}），用与之语义相近的已有 token 的嵌入来初始化。例如，\texttt{<|start\_header\_id|>} 可以用 “[” 或 “<” 等分隔符 token 的嵌入均值来初始化，因为它在功能上类似于一个开放分隔符。

\begin{corebox}[特殊 Token 设计的原则]
好的特殊 token 设计需要满足几个条件：
(1) \textbf{不与自然语言冲突}：
用 \texttt{<|...|>} 这样的格式是因为
它几乎不可能出现在正常文本中；
(2) \textbf{语义明确}：
模型和后处理代码都能明确判断每个 token 的含义；
(3) \textbf{向后兼容}：
添加新的特殊 token 不应破坏已训练的模型对旧 token 的理解。
这也是为什么添加特殊 token 通常伴随着一轮新的 SFT。
\end{corebox}

\section{分词的陷阱与挑战}

分词看似已经是一个「解决了的问题」，
但在实践中仍然存在许多令人头疼的陷阱。

\subsection{Glitch Token：分词器的「幻觉」}

2023 年，研究人员发现了一个诡异的现象：
当向 GPT 系列模型输入某些特定的 token
（如 “SolidGoldMagikarp”）时，
模型会表现出极不正常的行为：
重复输出无意义的文本、拒绝回答问题、
甚至产生与输入完全无关的内容。
% NEW REF: rumbelow2023glitch = Rumbelow, Jessica and watMEG, Matthew. SolidGoldMagikarp (plus, prompt generation). Lesswrong, 2023.

这些被称为 \textbf{glitch tokens}（故障 token）。
产生的原因是：
在 BPE 训练时，这些字符串在训练语料中频繁出现
（“SolidGoldMagikarp” 是一个 Reddit 用户名，
出现在大量被爬取的 Reddit 帖子中），
因此被编码为单独的 token。
但在模型的预训练数据中，
这些 token 可能很少或从未出现在有意义的上下文中：
它们的嵌入向量几乎没有被训练过，
停留在接近随机初始化的状态。

这个故事的教训是：
\textbf{分词器和模型的训练数据必须对齐}。
如果分词器的训练语料和模型的预训练语料差异过大，
就会产生大量「存在于词表中但模型不理解」的 token。

\subsection{防御 Glitch Token 的工程实践}

2024--2025 年，主流模型训练团队采用了几种方法来防范 glitch token 问题：

\textbf{词表修剪}。在分词器训练完成后，检查每个 token 在预训练语料中的实际出现频率。如果某个 token 在分词器训练语料中频繁出现（因此被合并为独立 token），但在预训练语料中几乎不出现，就将其从词表中移除。这种「分词器-预训练对齐」步骤已成为标准实践。

\textbf{嵌入热身}。在预训练的早期阶段，对词表中所有 token 的嵌入执行额外的正则化，确保每个 token 的嵌入至少被更新了最低次数。如果某个 token 在前 $N$ 步训练中从未出现过，可以主动用包含该 token 的合成样本来更新其嵌入。

\textbf{使用统一的训练语料}。最彻底的解决方案是让分词器的训练语料和模型的预训练语料完全相同（或至少是后者的子集）。DeepSeek-V3 明确采用了这一策略，分词器在与预训练相同的数据上训练，从源头消除了数据不对齐的问题。

\subsection{Token 嵌入的「死区」问题}

Glitch token 问题其实是一个更广泛现象的极端表现：\textbf{词表中大量 token 的嵌入训练不足}。

在一个典型的 128K 词表中，token 的频率分布遵循 Zipf 定律，最高频的 1000 个 token 占据了训练语料中超过 80\% 的 token 出现。而词表中排名在 100K 之后的 token，在整个训练过程中可能只出现了几千次。对于一个需要数十亿次梯度更新才能充分训练的嵌入向量来说，几千次出现是远远不够的。

这些训练不足的 token 形成了嵌入空间中的「死区」，它们的嵌入向量与其他 token 的嵌入没有形成有意义的关系。当模型在推理时遇到这些 token，其行为是不可预测的，可能产生幻觉、重复或完全无关的输出。

量化这个问题的一种方法是计算每个 token 嵌入的\textbf{有效训练次数}，即该 token 在训练过程中被梯度更新的总次数。经验法则是：一个 token 至少需要 $10^4$--$10^5$ 次有效训练才能形成稳定的嵌入。如果某个 token 的有效训练次数低于这个阈值，就应该被视为「欠训练」并在使用时加以注意。

\subsection{数学与数字的困境}

分词对模型的数学能力有着直观但容易被忽视的影响。
考虑 “12345” 这个数字：
在某些分词器下，它可能被切分为
“123” + “45”、“1” + “234” + “5”、
或者 “12345” 整体。
不同的切分方式意味着模型
「看到」的数字表示是不同的：
它甚至无法可靠地判断一个数字有几位。

这就是为什么许多现代分词器选择
\textbf{将每个数字字符独立为一个 token}：
“12345” $\to$ [“1”, “2”, “3”, “4”, “5”]。
虽然这增加了序列长度，
但让模型能够在字符级别理解数字的结构，
从而提升算术能力。
SentencePiece 的 \texttt{split\_digits=True}
选项正是为此设计的。

\subsubsection{数字位置嵌入：Abacus Embeddings}

即使将数字逐位切分，模型的算术能力仍然有限：
尤其在面对训练时未见过长度的数字时
（例如训练 20 位加法，推广到 100 位加法）。
McLeish 等人（2024）发现了问题的根源：
Transformer 的标准位置编码无法让模型
可靠地追踪每个数字在一个长数字中的\textbf{相对位置}
（个位、十位、百位……）。
% NEW REF: mcleish2024abacus = McLeish, Sean and Bansal, Arpit and Stein, Alex and Jain, Neel and Kirchenbauer, John and Bartoldson, Brian R and Kailkhura, Bhavya and Bhatele, Abhinav and Geiping, Jonas and Schwarzschild, Avi and Goldstein, Tom. Transformers Can Do Arithmetic with the Right Embeddings. In NeurIPS, 2024.

他们提出了 \textbf{Abacus Embeddings}：
为每个数字字符添加一个额外的嵌入，
编码该数字字符\textbf{相对于数字起始位置}的偏移量。
直觉上，这就像给每一位数字加上了
「我是个位」「我是十位」「我是百位」的标签。

结果令人印象深刻：
仅在 20 位数字上训练一天（单 GPU），
模型在 100 位加法上达到了 99\% 的准确率。
这个结果有深刻的启示：
\textbf{分词和嵌入的设计可以弥补（或加剧）
Transformer 在特定能力上的不足}。
数学能力不完全是「更多训练数据」的问题，
有时候只需要告诉模型数字的正确结构。

\subsubsection{数字比较的陷阱}

分词还导致了另一个反直觉的问题：数字大小比较。
许多 LLM 会错误地判断 “9.11 $>$ 9.9”：
因为 “.11” 和 “.9” 是不同的 token，
模型无法直接比较它们的数值大小。
这个问题的根源在于：
分词器将小数点后的数字视为文本片段而非数值，
“0.11” 中的 “.11” 可能被合并为一个 token，
而 “0.9” 中的 “.9” 是另一个 token：
模型看到的是两个完全不同的符号，
失去了数值比较的基础。

Kreitner 等人（2024）提出了 \textbf{xVal}，
一种直接将数字编码为单个 token 并附带浮点数值嵌入的方法，
在科学计算任务上显著提升了数字处理精度。
% NEW REF: kreitner2024xval = Kreitner, Linus and Hager, Paul and Mengedoht, Jonathan and Kaissis, Georgios and Rueckert, Daniel and Menten, Martin J. Efficient numeracy in language models through single-token number embeddings. arXiv preprint arXiv:2510.06824, 2024.

这些研究共同指向一个结论：
\textbf{BPE 是为自然语言设计的压缩算法，
不是为数学设计的}。
数字有独立于语言频率的内在结构（位值系统），
强行用统计频率来决定数字的切分方式，
从根本上就是一个范畴错误。

\subsubsection{科学计数法与浮点数的分词}

科学计算中的数字表示给分词带来了更多挑战。考虑 “3.14159e-10” 这个浮点数，它在分词器中可能被切为多种方式：“3” + “.” + “14” + “159” + “e” + “-” + “10”（7 个 token），或者 “3.14” + “159” + “e-” + “10”（4 个 token）。不同的切分方式意味着模型对这个数字的「理解」完全不同，在前一种切分下，“e” 看起来只是一个普通的字母，模型需要从上下文中推断它代表科学计数法。

更微妙的问题出现在负数和分数上。“-3/4” 中的 “-” 是负号还是减号？“3/4” 是分数还是日期格式？这些歧义在自然语言上下文中通常可以被消解，但分词器在做出切分决策时没有上下文信息，它只能基于训练语料中的频率统计来决定。

一些前沿研究开始探索\textbf{数字感知的分词}（number-aware tokenization）：识别文本中的数字，用专门的编码方式处理它们（如浮点嵌入或位值嵌入），然后将编码结果注入到 token 序列中。这种混合分词策略的工程复杂度较高，但在科学计算和金融领域可能带来显著的性能提升。

\subsection{代码中的空白问题}

Python 使用缩进来标记代码块层次。
一个常见的问题是：
分词器如何处理空格和制表符？

如果 4 个空格被编码为一个 token “\verb|    |”，
而 2 个空格被编码为另一个 token “\verb|  |”，
那么 8 个空格的缩进就是两个
“\verb|    |” token。
但如果分词器的合并规则不同，
8 个空格可能被切分为
“\verb|  |” + “\verb|    |” + “\verb|  |”
（3 个 token）：
这种不一致会让模型难以学习正确的缩进模式。

GPT-4 和 LLaMA 3 的分词器
为常见的空格序列（1 个到 16 个空格）
专门设置了独立的 token，
避免了这个问题。
同样，制表符和换行符也被设为独立 token。
这些看似微小的设计决策，
对模型的代码生成能力有着显著的影响。

一个定量的例子可以说明这一点。考虑以下 Python 代码片段：

\begin{lstlisting}[caption={Indentation-sensitive Python code}, language=Python]
def nested():
    for i in range(10):
        for j in range(10):
            if i == j:
                print(i)
\end{lstlisting}

如果分词器没有为空格序列设置专用 token，这段代码中 4 级缩进（16 个空格）可能被分为多种不一致的方式：有时是 4 个「4 空格」token，有时是 2 个「8 空格」token，有时是 1 个「16 空格」token。模型需要学习所有这些等价表示，这是一种不必要的学习负担。

而如果每种长度的空格序列都有专用 token，那么 4 级缩进\textbf{总是}被编码为 1 个「16 空格」token，确定性的分词使得模型可以建立「16 空格 = 4 级缩进」这样简洁的映射关系。SWE-bench 的研究表明，具有良好空白处理的分词器在代码修复任务上的准确率可以比差的分词器高出 3--5 个百分点。

\subsection{新编程语言的分词挑战}

随着 LLM 被用于越来越多的编程语言，
分词器面临着新的挑战：
尤其是那些在分词器训练语料中出现较少的新兴语言。

\textbf{Rust} 是一个典型的例子。
Rust 的语法中包含大量复合运算符
和在其他语言中少见的符号组合：
\texttt{::}（路径分隔符）、\texttt{->}（返回类型标记）、
\texttt{=>}（匹配分支）、\texttt{\&mut}（可变引用）、
\texttt{<'a>}（生命周期标注）。
在以 Python 和 JavaScript 为主训练的分词器中，
这些 Rust 特有的符号组合可能未被合并为独立 token，
而是被拆分为多个单字符 token：
增加了序列长度，也增加了模型学习 Rust 语法的难度。

以生命周期标注 \texttt{<'a, 'b>} 为例，
它可能被分为 7 个 token（\texttt{<, ', a, ,, ', b, >}），
而在 Rust 开发者眼中这是一个不可分割的语义单元。
更麻烦的是 Rust 的宏系统（如 \texttt{vec![1, 2, 3]}）：
\texttt{vec!} 中的感叹号在大多数语言中不出现在标识符后，
分词器会将其与标识符拆开，
丢失了「这是一个宏调用」的关键信息。

\textbf{Mojo}（Modular 推出的高性能 ML 语言）
面临更大的挑战：
它几乎不存在于 2024 年之前的训练语料中。
Mojo 的语法混合了 Python 和 Rust 的特征：
\texttt{fn}（函数定义）、\texttt{let}（不可变绑定）、
\texttt{var}（可变绑定）、\texttt{struct}、\texttt{trait}：
这些关键词虽然与 Rust 相同，
但 Mojo 独有的类型注解语法（如 \texttt{DType.float32}）
和 SIMD 操作（如 \texttt{SIMD[DType.float32, 4]}）
在训练语料中几乎不存在，
导致模型生成 Mojo 代码时错误率显著偏高。

\textbf{Zig}、\textbf{Gleam}、\textbf{Elixir} 等其他新兴语言面临类似的困境。
一个有趣的定量指标是「语言分词效率比」：
即同一功能的代码在不同语言中的 token 数比值。
以一个简单的 HTTP 服务器为例：
Python 版本约 50 token，
Rust 版本约 120 token（因为类型标注、生命周期等额外语法），
Mojo 版本约 150 token（因为 Mojo 特有的语法在词表中缺少合并规则）。
这意味着在相同的上下文窗口中，
模型能看到的 Mojo 代码比 Python 代码少约 3 倍：
这对代码理解和生成能力有直接影响。

\textbf{TypeScript} 是一个有趣的反例：
虽然它比 JavaScript 更冗长（因为类型标注），
但由于 TypeScript 在训练语料中的占比很高
（GitHub 上 TypeScript 仓库数量增长迅速），
分词器已经学会了为 TypeScript 特有的模式做高效的 token 合并。
例如，\texttt{interface\{} 和 \texttt{?: string} 在新一代分词器中
往往被合并为更紧凑的 token 序列。
这说明\textbf{分词效率本质上是由训练数据中该语言的占比决定的}：
更多数据 $\to$ 更好的 BPE 合并 $\to$ 更高的分词效率。

\begin{warnbox}[新语言的恶性循环]
新编程语言面临一个恶性循环：
(1) 训练语料中缺乏该语言的代码，
分词器无法为其优化 token 合并规则；
(2) 分词效率低导致模型生成该语言代码的质量差；
(3) 质量差导致开发者不使用 LLM 辅助编写该语言代码，
进一步减少了高质量训练数据的来源。
打破这个循环的方法通常是
在微调阶段大量注入特定语言的高质量代码数据：
但这需要分词器的词表中
已经包含该语言的常见 token，
否则效果仍然有限。
\end{warnbox}

\section{分词器的调试与诊断}

在实际工作中，分词器的 bug 往往是最难发现的，因为它们不会产生明确的错误消息，只会导致模型性能的「静默下降」。以下是几种常见的分词问题及其诊断方法。

\subsection{可视化分词结果}

调试分词器的第一步是\textbf{可视化分词结果}。将文本的分词结果以彩色标注的方式展示，可以直观地发现不合理的切分。例如：

\begin{lstlisting}[caption={Visualizing tokenization with tiktoken}, language=Python]
import tiktoken

enc = tiktoken.encoding_for_model("gpt-4")
text = "DeepSeek-R1 achieves state-of-the-art reasoning."
tokens = enc.encode(text)

# Visualize each token
for t in tokens:
    decoded = enc.decode([t])
    print(f"[{decoded}]", end=" ")
# Output: [Deep] [Seek] [-R] [1] [ achieves] [ state] [-of] [-the] [-art]
#         [ reasoning] [.]
\end{lstlisting}

从这个输出中可以观察到几个特点：(1) “DeepSeek” 被拆为 “Deep” + “Seek”（因为 “DeepSeek” 在训练时可能不是高频词）；(2) 连字符 “-” 被和前后的词合并而非独立存在；(3) 空格被合并到后面的词中（如 “ achieves” 包含了前导空格）。

\subsection{频率分析}

一种更系统的诊断方法是对训练语料中所有 token 的频率做统计分析。如果发现大量 token 的出现频率极低（例如低于 1000 次），这些 token 可能是 glitch token 的候选者。可以用以下指标来评估词表的「健康度」：

\begin{itemize}[noitemsep]
  \item \textbf{活跃 token 比例}：在训练语料中至少出现 $10^4$ 次的 token 占词表的百分比。理想值 $>$ 90\%。
  \item \textbf{长尾 token 比例}：出现次数不到 100 的 token 占词表的百分比。理想值 $<$ 5\%。
  \item \textbf{最高/最低频率比}：词表中最高频 token 和最低频 token 出现次数的比值。如果这个比值超过 $10^6$，说明词表中存在大量「死」token。
\end{itemize}

\subsection{跨语言 fertility 对比}

对于多语言模型，定期测量不同语言的 fertility 是发现分词偏见的有效手段。如果某种语言的 fertility 显著高于其他语言（例如超过 2.0），说明分词器对该语言的支持不足，可能需要扩充词表或调整训练数据的语言比例。

\section{分词对下游性能的影响}

分词看似只是预处理，但它的影响贯穿整个模型的生命周期。本节系统性地分析分词选择如何影响下游任务的表现。

\subsection{分词器作为关键超参数}

长期以来，分词器在 LLM 训练中被视为一个「设置后就忘掉」的组件，研究者把绝大部分精力投入到架构设计、训练策略和数据工程中，对分词器只是选择一个「够用的」方案就继续往前走。这种态度在 2024 年开始改变。

2024 年的多项研究~\cite{goldman2024tokenizer} 表明，\textbf{分词器的选择对模型性能的影响可以等价于改变 30\%--50\% 的训练数据量}。这个发现颠覆了此前「分词器不重要，只要能用就行」的普遍认知。

具体来说，研究者在相同的模型架构和训练数据上，仅改变分词器（BPE vs Unigram，32K vs 64K vs 128K 词表），观察到以下现象：

\textbf{通用语言理解}（MMLU、HellaSwag 等）对分词器不太敏感，性能差异在 1--2\% 以内。这是因为这些任务主要测试高层语义理解，分词粒度的影响被 Transformer 的多层抽象吸收了。

\textbf{数学和推理}（GSM8K、MATH）对分词器非常敏感，性能差异可达 5--10\%。这与前文讨论的数字分词问题直接相关：不同的数字切分方式会显著影响模型的算术能力。

\textbf{多语言任务}对分词器极度敏感，性能差异可达 15--20\%。一个对特定语言 fertility 很高的分词器，等于在该语言上「浪费」了大量的上下文窗口和计算资源。

\textbf{代码生成}对分词器中度敏感，性能差异在 3--8\%。关键因素是分词器能否正确处理缩进、运算符和编程语言特有的符号。

\subsection{分词器选择的实际影响案例}

为了让分词器的影响更加具体，让我们看几个真实的案例。

\textbf{案例一：LLaMA 1 的中文能力瓶颈}。
LLaMA 1 使用了 32K 的 SentencePiece 词表，以英语为主训练。
当中文 AI 社区尝试在 LLaMA 1 上微调中文能力时，
发现了一个根本性的瓶颈：
即使使用大量高质量的中文 SFT 数据，
模型的中文理解和生成能力仍然远低于英语。
根因分析发现：分词器将中文文本的 fertility 推高至 2.5，
使得模型在处理同等长度的中文内容时，
有效上下文窗口只有英语的 40\%。
这不是微调可以解决的问题，\textbf{分词器的限制是架构层面的}。
Chinese-LLaMA~\cite{cui2024llama} 通过扩充词表才部分解决了这个问题。

\textbf{案例二：GPT-2 的 glitch token 危机}。
GPT-2 的分词器在 Reddit 数据上训练，
导致一些 Reddit 特有的用户名（如 “SolidGoldMagikarp”）
被编码为独立的 token。
当这些 token 出现在模型输入中时，
由于它们的嵌入几乎未被训练，
模型会产生完全不可预测的输出：
包括重复输出、胡言乱语和拒绝服务。
这个案例促使后来的模型团队
更加重视分词器和预训练数据的对齐。

\textbf{案例三：代码模型的缩进 bug}。
某些早期的代码模型在生成 Python 代码时，
会偶尔产生混合制表符和空格的缩进：
这在 Python 中是语法错误。
根因是分词器没有为制表符和空格序列
设置独立的 token，
导致模型在生成缩进时混淆了不同的空白表示。
这个 bug 在 GPT-4 和 LLaMA 3 中通过
显式的空白 token 设计得到了解决。

\subsection{衡量分词器质量的量化指标}

如何系统地评估一个分词器的质量？以下几个指标在实践中最为常用：

\textbf{Fertility（生育率）}，即平均每个词被切分为多少个 token。fertility 越接近 1 越好。这是最直接的效率指标，可以针对不同语言单独计算：

\begin{equation}
  \text{fertility}(\text{language}) = \frac{\text{总 token 数}}{\text{总词数（空格分隔）}}
\end{equation}

对于中文等无空格语言，通常以字符为基准：$\text{fertility} = \text{总 token 数} / \text{总字符数}$。

\textbf{压缩率}，即分词后的 token 序列长度与原始字节长度的比值。压缩率越低越好（更少的 token 表示相同内容）。这与 fertility 相关但不同，fertility 以「词」为基准，压缩率以「字节」为基准：

\begin{equation}
  \text{compression ratio} = \frac{\text{token 序列长度}}{\text{原始 UTF-8 字节长度}}
\end{equation}

\textbf{词表覆盖率}，即训练语料中不同 token 被使用的比例。如果词表中 20\% 的 token 在训练过程中从未出现，这些 token 就是浪费的词表空间。理想情况下，每个 token 至少在训练过程中出现 $10^4$ 次以上。

\textbf{分词一致性}，即语义相似的词是否产生结构相似的分词结果。例如，“playing”、“played”、“plays” 是否都共享 “play” 这个子词？好的分词器应该在形态学变体之间保持一致的子词结构。

\textbf{续写连贯性}，在分词边界处，模型是否能自然地续写？如果 “transformers” 被切为 “trans” + “formers”，模型在看到 “trans” 后是否能合理地预测 “formers”？可以通过在分词边界处测量模型的 perplexity 来量化这个指标，在不合理的分词边界处，perplexity 通常会异常偏高。

\subsection{一个完整的分词器评估案例}

让我们用一个具体的例子来说明如何系统地评估分词器。假设我们要在以下四个维度上比较 LLaMA 1（32K 词表）和 Qwen3（151K 词表）的分词器：

\textbf{英语效率}。在一段 10,000 词的英文新闻文本上：
LLaMA 1 产生约 13,500 个 token（fertility 1.35），
Qwen3 产生约 12,800 个 token（fertility 1.28）。
差异不大，因为两者都对英语有充分的优化。

\textbf{中文效率}。在一段等长的中文新闻文本上（约 15,000 个汉字）：
LLaMA 1 产生约 38,000 个 token（fertility 2.53/字符），
Qwen3 产生约 18,500 个 token（fertility 1.23/字符）。
差异巨大，Qwen3 的中文效率是 LLaMA 1 的约 2 倍。

\textbf{代码效率}。在一段 5,000 行的 Python 代码上：
LLaMA 1 产生约 85,000 个 token，
Qwen3 产生约 72,000 个 token。
Qwen3 的代码效率更高，主要得益于其词表中包含了更多的编程相关 token（如完整的关键词和常见的 API 名称）。

\textbf{数学效率}。在一段包含大量数学公式的 LaTeX 文本上：
两者的 fertility 差异不大（因为数学符号在两个词表中都以单字符 token 存在），但 Qwen3 对 LaTeX 命令（如 \texttt{\textbackslash frac}、\texttt{\textbackslash sum}）有更好的合并，将常见的 LaTeX 命令合并为独立 token，减少了序列长度。

这个评估案例说明了一个重要的实践原则：\textbf{分词器的评估必须覆盖目标应用的所有语言和领域}。一个在英语上表现优异的分词器可能在中文上灾难性地低效，而反之亦然。

\subsection{Token 经济学：分词效率的商业影响}

在 API 服务时代，分词效率有直接的商业影响。以下计算说明了这一点。

假设一个企业每月通过 API 处理 10 亿个中文字符。使用 LLaMA 1 分词器（fertility 2.5），这相当于 25 亿 token；使用 Qwen3 分词器（fertility 1.2），这相当于 12 亿 token。如果 API 定价为每百万 token 2 美元：

\begin{itemize}[noitemsep]
  \item LLaMA 1 分词器的月成本：$25 \times 10^3 \times 2 = \$50{,}000$
  \item Qwen3 分词器的月成本：$12 \times 10^3 \times 2 = \$24{,}000$
  \item 年节约额：$\$312{,}000$
\end{itemize}

仅仅因为分词器的差异，企业每年就可以节省超过 30 万美元。这还没有考虑分词效率带来的上下文窗口收益，使用 Qwen3 分词器时，同样的上下文窗口可以容纳约 2 倍的中文内容，这意味着模型在更长的文档上的理解能力更强，从而间接提升了服务质量。

这种「token 经济学」是中文 AI 社区强烈推动大词表分词器的直接驱动力之一。Qwen、DeepSeek、GLM 等中国模型无一例外地使用了 100K+ 的大词表，因为对于中文用户来说，分词效率的提升等同于\textbf{实质性的成本降低和能力提升}。

\subsection{分词与推理链的关系}

一个鲜为人知但极其重要的观察是：分词粒度会影响模型的\textbf{推理深度}。

Transformer 模型的每一层对每个 token 执行一次注意力计算和前馈变换。如果一个概念被编码为 1 个 token，它经历 $L$ 次变换（$L$ 是层数）。如果同样的概念被编码为 $k$ 个 token，每个 token 经历 $L$ 次变换，并且这些 token 之间通过注意力机制互相影响，这意味着模型有效地拥有了更多的「计算步骤」来处理这个概念。

这个观察有一个反直觉的推论：\textbf{更长的 token 序列可能给模型更多的「思考空间」}。这部分解释了为什么 Chain-of-Thought prompting 如此有效，通过让模型显式地写出中间步骤，模型获得了额外的 token 位置来执行中间计算~\cite{wei2022cot}。从这个角度看，分词器的粒度选择不仅是效率问题，更是\textbf{计算深度}问题。

这并不意味着更细的分词总是更好，序列长度增加也会导致注意力计算成本的二次增长，以及远距离依赖学习的困难。最优的分词粒度是效率和能力之间的平衡，而这个平衡点取决于具体任务的计算需求。

\subsection{分词对安全性的隐性影响}

分词器的设计还会影响模型的安全行为，这一点鲜为人知但后果严重。

\textbf{对抗性 token 拼接}。攻击者可以利用分词器的特性来绕过安全过滤。例如，如果「如何制造\_\_\_\_」被安全过滤器拦截，攻击者可能使用不同的空格或 Unicode 字符来改变分词结果，使得安全过滤器无法匹配到危险模式，但模型仍然能理解语义。某些 Unicode 同形字（如西里尔字母 “а” 看起来和拉丁字母 “a” 相同，但 UTF-8 编码不同）可以产生完全不同的 token 序列，使基于 token 匹配的安全过滤器失效。

\textbf{跨语言安全漏洞}。安全训练（RLHF/DPO）通常以英语为主。但如果分词器对某种语言的支持不佳（高 fertility），模型在该语言上的安全行为可能不如英语可靠。研究表明~\cite{deng2024multilingual}，某些模型在收到低资源语言的有害请求时，拒绝率显著低于英语，这部分原因在于分词器将低资源语言的有害内容编码为了与英语安全训练数据完全不同的 token 序列，使得安全对齐无法泛化。

\textbf{Prompt injection 的分词维度}。Prompt injection 攻击的一些变体利用了分词器处理特殊字符的方式。例如，某些不可见的 Unicode 控制字符（如零宽度空格 U+200B）可能被分词器忽略或映射为特殊 token，但在模型的理解中产生了「分隔」效果，攻击者可以利用这种分隔来「隐藏」恶意指令。

\subsection{词表大小的实践建议}

综合以上分析，以下是 2025--2026 年分词器设计的最佳实践：

\begin{keybox}[分词器设计的最佳实践]
\begin{itemize}[noitemsep]
  \item \textbf{模型 $\geq$ 7B}：词表大小 64K--150K 是合理范围。过小（32K）会严重影响多语言和代码性能，过大（256K+）在小模型上浪费参数预算。
  \item \textbf{模型 $\geq$ 70B}：词表大小 128K--262K。嵌入层开销可忽略，大词表的多语言和代码收益值得追求。
  \item \textbf{强制 split digits}：将每个数字字符独立为一个 token。这是提升数学能力的最简单有效手段。
  \item \textbf{平衡语言采样}：分词器训练数据中，目标语言的比例应不低于 10\%。如果目标包括中文，CJK 数据比例应至少 20\%。
  \item \textbf{代码空白 token}：为 1--16 个空格的序列设置专用 token。制表符和换行符也应独立。
  \item \textbf{与预训练数据对齐}：分词器训练数据应是预训练数据的子集或近似分布，避免 glitch token 问题。
\end{itemize}
\end{keybox}

\section{前沿探索：超越传统分词}

\subsection{Tiktoken：速度的极致}

2022 年，OpenAI 开源了 tiktoken：
一个用 Rust 编写的高性能 BPE 分词器。
tiktoken 的设计哲学是\textbf{推理速度优先}：
它使用正则表达式进行预分词（pre-tokenization），
将文本先切分为较大的块，
然后在每个块内做 BPE 匹配。
这种设计使得 tiktoken 的编码速度
比 Python 实现快 3--6 倍。

GPT-4 和 GPT-4o 都使用 tiktoken。
LLaMA 3 也从 SentencePiece 切换到了
类似 tiktoken 的 BPE 实现。

tiktoken 和 SentencePiece 的核心区别不在算法层面（都支持 BPE），而在\textbf{预分词策略}上。SentencePiece 将输入视为连续的字节流，不做预分词，BPE 的合并可以跨越「词」边界。tiktoken 则先用正则表达式将文本切分为独立的片段（通常对应于「词」或标点符号），然后在每个片段内部独立地做 BPE，合并永远不会跨越片段边界。

这个看似微小的差异有实际影响。在 SentencePiece 中，“New York” 中的空格可能被合并到前一个或后一个词中，形成 “New\_York” 这样的 token，两个词被合并为一个 token。在 tiktoken 中，”New” 和 “York” 被预分词步骤隔离在不同的片段中，永远不会被合并。哪种更好？没有定论，合并可以减少 token 数但可能丢失词边界信息，不合并保持了词的独立性但增加了 token 数。

\subsection{分词器的吞吐量基准}

在生产环境中，分词器的速度也是一个重要考量。以下是几种主流分词器实现在标准基准上的编码速度对比：

\begin{center}\small
\renewcommand{\arraystretch}{1.3}
\begin{tabular}{lrl}
\toprule
\rowcolor{figLightGray}
\textbf{实现} & \textbf{编码速度 (MB/s)} & \textbf{语言} \\
\midrule
SentencePiece (C++) & 25--40  & C++ \\
tiktoken (Rust)     & 80--150 & Rust + Python bindings \\
HuggingFace tokenizers (Rust) & 100--200 & Rust + Python bindings \\
Python BPE (参考实现) & 2--5 & Pure Python \\
\bottomrule
\end{tabular}
\end{center}

Rust 实现的分词器比纯 Python 实现快 20--50 倍，这在处理 TB 级训练语料或高吞吐量推理服务时是决定性的差异。HuggingFace 的 \texttt{tokenizers} 库因其高性能和广泛的模型支持，已成为分词器的事实标准实现。

\subsection{无分词器模型：字节级别的终极方案？}

一个更激进的想法是完全抛弃分词器，
直接在字节级别训练模型。

\subsubsection{ByT5：最早的纯字节模型}

Google 在 2022 年发表的 ByT5~\cite{xue2022byt5} 是最早的纯字节级 Transformer 模型之一。ByT5 直接处理 UTF-8 字节序列，不使用任何分词器，词表只有 256 个字节 token 加上少量特殊 token。

ByT5 的实验结果揭示了字节级模型的一个有趣的权衡：\textbf{在字符级任务上（拼写纠错、音素转写、形态学分析），ByT5 显著优于基于 token 的 mT5；但在语义理解任务上，两者性能相当或 ByT5 稍劣}。这个结果暗示了一个深层原因：字节级模型在底层处理上更精确（因为它「看到」了每一个字符），但在高层语义抽象上可能需要更多的层来弥补信息粒度过细的问题。

ByT5 的致命弱点是\textbf{序列长度爆炸}。同样一段文本，字节表示的序列长度大约是 BPE 表示的 4--6 倍。对于 Transformer 的 $O(n^2)$ 注意力机制来说，这意味着计算量增加了 16--36 倍。ByT5 通过增加 encoder 层数（12 层 $\to$ 36 层）和减少 decoder 层数来缓解这个问题，但总体效率仍然远低于基于 token 的模型。

\subsubsection{MegaByte：分层字节处理}

Meta 的 MegaByte（2023 年）提出了更实用的方案：
它将输入视为原始字节序列，
用一个大模型处理字节块，
用一个小模型在块内生成单个字节。
% NEW REF: yu2023megabyte = Yu, Lili and Simig, Daniel and Flaherty, Colin and Aghajanyan, Armen and Zettlemoyer, Luke and Lewis, Mike. MEGABYTE: Predicting Million-Byte Sequences with Multiscale Transformers. In NeurIPS, 2023.

MegaByte 的分层设计解决了 ByT5 的效率问题：大模型（Global Model）在块级别操作，每个块包含 $P$ 个字节（通常 $P = 4$ 或 $P = 8$），因此大模型的序列长度只有原始字节序列的 $1/P$。小模型（Local Model）则在每个块内部逐字节生成，但由于块很短，计算量很小。

这种设计的问题在于\textbf{块边界是固定的}，不考虑文本内容。一个词可能被分到两个不同的块中（例如 “th\textbar e” 被分在相邻两个块），导致大模型无法完整地看到词内结构。这个问题在 BLT 中得到了优雅的解决。

字节级模型的潜在优势非常诱人：
\begin{itemize}[noitemsep]
  \item 完全消除分词器，简化整个训练和部署流程。
  \item 真正的语言无关，任何语言、任何编码都是字节。
  \item 消除分词导致的信息损失和偏见。
\end{itemize}

\subsubsection{BLT：动态分配计算的字节级模型}

2024 年 12 月，Meta 发表了 \textbf{Byte Latent Transformer}（BLT），
标志着字节级模型在实用性上迈出了关键一步。
BLT 的核心创新是\textbf{动态 patch 分割}：
它不使用固定大小的字节块（如 MegaByte 的做法），
而是根据\textbf{下一个字节的熵}来决定 patch 的边界。
% NEW REF: pagnoni2024blt = Pagnoni, Artidoro and Pasunuru, Ram and Rodriguez, Pedro and Nguyen, John and Muller, Benjamin and Li, Margaret and Zhou, Chunting and Yu, Lili and Weston, Jason and Zettlemoyer, Luke and Ghosh, Gargi and Lewis, Mike and Holtzman, Ari and Iyer, Srinivasan. Byte Latent Transformer: Patches Scale Better Than Tokens. arXiv preprint arXiv:2412.09871, 2024.

直觉上，这种设计非常自然：
\begin{itemize}[noitemsep]
  \item 当文本可预测时（例如常见的英文单词中间），
        熵低，patch 变长，用更少的计算处理简单内容。
  \item 当文本不可预测时（例如代码中的变量名、罕见词），
        熵高，patch 变短——分配更多计算给复杂内容。
\end{itemize}

BLT 的架构由三个组件构成：
一个轻量的\textbf{字节编码器}将字节序列编码为 patch 表示，
一个大型的\textbf{潜在 Transformer} 在 patch 级别进行推理，
一个轻量的\textbf{字节解码器}将 patch 表示还原为字节预测。
重型计算（注意力、FFN）只在 patch 级别发生，
而 patch 的平均长度约为 4--6 个字节：
接近 BPE token 的平均长度，
因此计算成本可以与基于 token 的模型竞争。

Meta 的实验结果令人瞩目：
在 8B 参数、4T 训练字节的规模上，
BLT 首次在 FLOP 控制的对比中
\textbf{匹配了基于 token 的 LLaMA 模型}的性能，
同时在推理效率上有所提升
（因为简单文本使用更长的 patch，减少了前向传播次数）。
更重要的是，BLT 展现了更好的\textbf{长尾泛化能力}：
在拼写、音素等字符级任务上
显著优于基于 BPE 的模型。

\begin{corebox}[BLT 的意义]
BLT 证明了字节级模型的可行性不再是理论上的。
它的关键洞察是：
\textbf{计算资源应该根据内容的复杂度动态分配，
而不是像 BPE 那样用固定的词表来决定}。
这与人类阅读的模式惊人地一致：
我们扫过熟悉的词很快，
在遇到生僻词或复杂概念时则放慢速度仔细理解。
截至 2026 年初，BLT 的代码已开源，
但尚未有 frontier 级别的模型采用纯字节级架构。
这个方向值得密切关注。
\end{corebox}

\subsubsection{ByteFlow Net：第二代字节级模型}

2026 年 3 月 Amazon 发布的 \textbf{ByteFlow Net}~\cite{amazon2026byteflow}
把字节级建模推向了真正端到端的形态。
BLT 的 patch 边界由独立训练的字节级熵模型给出，
是一种「先训熵模型、再用全局阈值切分」的两阶段流水线；
ByteFlow Net 则取消了这个外部依赖，
直接在原始字节流上学习语义有意义的分块边界。

ByteFlow Net 的核心创新是
用\textbf{潜在表示的 coding rate}
（描述一组向量在嵌入空间中所占的「编码体积」）
作为分块准则：
当连续若干字节的潜在表示已经
能被一个紧凑的子空间充分覆盖时，
就在该位置切分，开始新的 patch。
为了保持静态计算图，
ByteFlow Net 在每一步用 Top-K 选择
而非阈值比较来决定边界——
这避免了动态形状带来的工程复杂度，
让字节级 patch 机制可以直接复用现有的 Transformer 编译器栈。

BLT 与 ByteFlow Net 的关键差异在于：
BLT 的熵模型是\textbf{静态}的（训练完即冻结），
全局阈值在数据分布漂移时未必最优；
ByteFlow Net 的 coding rate 信号
完全由当前模型的内部表示给出，
对输入自适应、对训练动态自适应。
在公开的 FLOP 控制对比中，
ByteFlow Net 同时超过了
基于 BPE 的标准 Transformer
和 BLT 风格的两阶段字节级 baseline，
且差距在长尾、罕见模式上最显著。

ByteFlow Net 的意义并非又一次单点性能提升，
而是把字节级模型的工程门槛压低到了
与基于 token 模型可比的水平：
没有独立的熵模型要训练、
没有需要调的全局阈值、
没有需要兼容动态形状的特殊编译器。
若这条路线得到验证，
\textbf{字节级模型可能在 2026--2027 年间
首次具备替代 BPE 的实用条件}。

\subsubsection{字节级模型的前景评估}

截至 2026 年 3 月，字节级模型仍处于「概念验证」到「早期实用」的过渡阶段。以下几个因素将决定它们能否在未来取代 BPE：

\textbf{训练效率的差距}。尽管 BLT 在 FLOP 对比中匹配了 LLaMA，但在实际的 wall-clock 训练时间中仍有差距。字节级模型的序列更长，对 GPU 内存和通信带宽的要求更高。在分布式训练中，这个差距可能被进一步放大。字节级模型在训练时需要处理约 4--6 倍长的序列，即使使用 patch 机制将注意力计算限制在 patch 级别，字节编码器和解码器的额外计算仍然是一个不可忽视的开销。

\textbf{长上下文的天然优势}。随着上下文窗口从 128K 扩展到 1M+ token，字节级模型的一个潜在优势浮现了：它们不需要在扩展上下文时重新训练分词器。基于 token 的模型在处理超长上下文时，分词器的固定词表可能成为瓶颈（某些极长的文档中包含的模式在训练词表时未被覆盖）。字节级模型天然地避免了这个问题。

\textbf{多模态统一的契合性}。字节级表示与多模态 token 有天然的兼容性，图像、音频、视频的原始表示都是字节流。一个统一的字节级模型可以直接处理所有模态的原始数据，无需为每种模态设计专门的分词器。这个方向与 LLaMA 4 和 Gemini 的多模态统一趋势高度吻合。

\textbf{字节级模型 vs BPE 的定量对比}。以下表格总结了截至 2026 年初，字节级模型和 BPE 模型在关键维度上的对比：

\begin{center}\small
\renewcommand{\arraystretch}{1.3}
\begin{tabular}{lcc}
\toprule
\rowcolor{figLightGray}
\textbf{维度} & \textbf{BPE 模型} & \textbf{字节级模型} \\
\midrule
词表大小       & 32K--262K     & 256 (+特殊 token) \\
序列长度倍数   & 1$\times$     & 4--6$\times$ \\
字符级任务     & 较弱           & 显著更强 \\
语义理解       & 强             & 接近（BLT 8B） \\
训练效率       & 基准           & 0.7--0.9$\times$ \\
多语言公平性   & 取决于词表     & 天然公平 \\
分词器维护     & 需要           & 不需要 \\
代表模型       & GPT/LLaMA/Qwen & ByT5/MegaByte/BLT \\
\bottomrule
\end{tabular}
\end{center}

这个对比显示了字节级模型的核心权衡：以训练效率的小幅下降换取了分词器的完全消除和多语言公平性。随着硬件效率的持续提升（特别是注意力机制的硬件加速），训练效率的差距可能进一步缩小，使字节级模型变得越来越有吸引力。

\subsection{多模态分词：图像也是 Token}

随着多模态大模型（如 GPT-4V、Gemini）的兴起，
一个自然的问题是：
能不能把图像也「分词」成一组离散 token，
与文本 token 统一处理？

答案是肯定的。
VQ-VAE（Vector Quantized Variational Autoencoder）
和 VQGAN 等方法可以将图像编码为
一组离散的码本索引：
每个索引就是一个「图像 token」。
例如，一张 256$\times$256 的图像
可以被编码为 $16 \times 16 = 256$ 个图像 token。
这些 token 与文本 token 共享同一个 Transformer，
实现了真正的多模态统一建模。

这种方法的代表包括：
Google 的 Parti 和 Gemini、
Meta 的 Chameleon。
它们展示了一个激动人心的可能性：
\textbf{所有模态（文本、图像、音频、视频）
都可以被「分词」为离散 token，
然后用同一个 Transformer 进行统一建模}。
分词不再只是 NLP 的预处理步骤，
而是多模态 AI 的基础设施。

\subsubsection{Emu3：纯 next-token prediction 的多模态统一}

2024 年，北京智源人工智能研究院（BAAI）发布了 \textbf{Emu3}，
这是多模态分词领域的一个里程碑式工作。
Emu3 证明了一个强有力的命题：
\textbf{仅靠 next-token prediction，
不需要扩散模型或组合架构，
就可以同时实现多模态感知和生成}。
% NEW REF: wang2024emu3 = Wang, Xinlong and Zhang, Xiaosong and Luo, Zhengxiong and Sun, Quan and Cui, Yufeng and others. Emu3: Next-Token Prediction is All You Need. Nature, 2025.

Emu3 的方法概念上非常简洁：
将图像、文本和视频统一编码为离散 token 序列，
然后用一个标准的 Transformer decoder 做 next-token prediction。
图像通过改进的 VQVAE 编码为 $32 \times 32 = 1024$ 个视觉 token，
视频则被编码为帧序列中的视觉 token 流。
在推理时，生成图像就是自回归地生成视觉 token，
然后用 VQVAE decoder 解码回像素。

Emu3 在图像生成上超越了 SDXL，
在图像理解上与 LLaVA-1.6 持平：
这是第一次有纯自回归模型在两个方向上同时达到领先水平。
这一结果发表在 \textit{Nature}（2025）上，
标志着多模态分词从实验性概念走向了成熟验证。

\subsubsection{UniTok：统一视觉分词器}

多模态分词面临的一个技术挑战是
\textbf{生成和理解对视觉 token 的需求不同}。
图像生成需要高保真度的像素重建
（偏向 VQVAE 式的离散编码），
而图像理解需要高级语义特征
（偏向 CLIP 式的连续编码）。

ByteDance 和香港大学的 \textbf{UniTok}（2025）
尝试在一个统一的分词器中同时满足两种需求。
UniTok 的核心创新是在训练视觉分词器时
同时优化重建损失和语义对齐损失，
使得同一组离散 token 既能重建图像，
也能承载语义信息。
这个方向代表了多模态分词的未来趋势：
不是为每个任务设计专门的分词器，
而是用一个统一的分词器覆盖所有需求。
% NEW REF: ma2025unitok = Ma, Chuofan and Jiang, Yi and Wu, Junfeng and Yang, Jihan and Yu, Xin and Yuan, Zehuan and Peng, Bingyue and Qi, Xiaojuan. UniTok: a Unified Tokenizer for Visual Generation and Understanding. In NeurIPS, 2025.

\subsubsection{LongCat-Next：DiNA 与 dNaViT 的统一离散空间}

2026 年 4 月美团开源的 \textbf{LongCat-Next}~\cite{meituan2026longcatnext}
把统一 tokenizer 路线推进到了一个新阶段。
它建立在 \textbf{DiNA}（Discrete Native Autoregressive）框架之上：
文本、视觉、音频被映射到同一个离散 token 空间，
所有模态共用同一个 next-token prediction 目标进行训练。

LongCat-Next 的核心组件是 \textbf{dNaViT}——
一个支持任意分辨率的视觉 tokenizer 与 de-tokenizer。
传统 ViT 系列的视觉编码器
要么把输入压成固定 patch 网格、
要么把图像缩放到训练分辨率，
两种做法都会对长宽比和细节分布造成失真。
dNaViT 直接处理原生分辨率的视觉信号，
将其编码为\textbf{分层离散 token}：
全局结构用稀疏的低层 token 表示，
细节区域用密集的高层 token 表示，
然后由同一个 Transformer 在 token 级别解释它们。

LongCat-Next 试图解决统一 tokenizer 此前的两个长期痛点。
第一是\textbf{离散视觉建模在理解任务上的性能上限}。
Chameleon 等早期工作虽然实现了文本与图像的统一离散建模，
但在视觉理解基准上长期落后于 CLIP/SigLIP 风格的连续编码器；
dNaViT 的分层离散设计在不放弃统一空间的前提下
重新拉近了与连续编码器的差距。
第二是\textbf{理解与生成之间的内在冲突}。
理解偏好语义抽象的 token，
生成偏好像素重建保真的 token，
UniTok 通过双目标训练做了一种妥协；
LongCat-Next 则通过分层 token 自然地把两种需求分配到不同层级。

把这条路线放到历史脉络中看：
Chameleon（2024）证明了文本与图像可以共享离散空间，
AnyGPT（2024）把这个空间扩展到了音频，
LongCat-Next 是这条路线在 2026 年最重要的延续，
它标志着「统一 tokenizer」
从概念验证走向了工程上可复现的开源基线。

\subsubsection{音频和视频的分词}

音频和视频的分词也在快速发展。
音频方面，\textbf{SoundStream}（Google）和
\textbf{EnCodec}（Meta）使用残差向量量化
（Residual Vector Quantization, RVQ）
将音频波形编码为多层离散 token。
Gemma 3n（2025）直接以 6.25 tokens/秒 的速率
处理音频输入：
这意味着 1 分钟的音频仅需要 375 个 token，
远低于将音频转录为文本的 token 开销。

视频则是最具挑战性的模态。
一段 10 秒、30fps 的 256$\times$256 视频，
如果每帧编码为 256 个 token，
总共需要 $300 \times 256 = 76{,}800$ 个 token：
即使对于 10M 上下文窗口的 LLaMA 4 Scout
也是一个不小的负担。
视频分词的关键是\textbf{时间冗余压缩}：
相邻帧之间的变化通常很小，
只需要编码变化的部分，
类似于视频编解码器中 I 帧和 P 帧的关系。

一个值得关注的方向是\textbf{自适应分辨率的视觉分词}。
不同的图像区域包含不同密度的信息：
一张照片中天空占了大半（信息密度低，用少量 token 即可），
而物体的细节区域信息密度高，需要更多 token。
SPHINX（2024）和 LLaVA-Next 等工作开始探索
根据图像内容自适应地分配视觉 token 数量：
信息密度高的区域分配更多 token，
信息密度低的区域分配更少 token。
这种做法与 BLT 在文本上的动态 patch 策略异曲同工：
\textbf{计算资源应该根据内容复杂度动态分配}。

\subsubsection{音乐 tokenizer 的个性化：Suno v5.5}

音乐生成长期被视作多模态 tokenizer 的一个独立分支。
2026 年 3 月 26 日发布的 \textbf{Suno v5.5}~\cite{suno2026v55}
把这个分支的产品形态推到了一个新阶段。
v5.5 的三项核心更新都直接作用在「音乐 token 的个性化」上。
\textbf{Voices} 允许 Pro/Premier 用户上传自己的声音作为生成条件，
为防止滥用，
开放使用前需要朗读一段验证短语。
\textbf{Custom Models} 让每位 Premier 用户最多保留 3 个个性化的 v5.5 实例，
每个实例都可以在自己的训练样本上微调音乐 tokenizer 与生成头。
\textbf{My Taste} 则是一个偏好学习层，
通过用户对生成结果的隐性反馈
持续调整解码策略。

时长约束同样在演进。
单次生成的最长时长在 v5 上是 8 分钟，
Premier 订阅可通过链式扩展把单首作品延长到 12 分钟以上。
版权方面，
Warner Music 已在 2025 年 11 月与 Suno 和解，
Sony 与 Universal 仍在诉讼中——
这意味着「音乐 tokenizer 的训练数据是否合法」
仍是一个未决问题。

把 Suno v5.5 放到本章脉络中看，
它示范了「tokenizer 个性化」的工程含义：
当一个模态的离散 token 空间
已经稳定到可以做产品化时，
下一阶段的差异化竞争
不再发生在 tokenizer 本身的算法层，
而发生在「让 tokenizer 学会用户偏好」的接口层。
这条规律对未来的图像、视频生成产品同样适用。

\subsubsection{ByteDance Omni 与 Context Unrolling：五模态联合预训练}

2026 年 4 月 23 日 ByteDance 发布的 \textbf{Omni}~\cite{bytedance2026omni}
把「统一离散空间」的范畴推到了五种模态：
text、image、video、3D-geometry、hidden-representation。
最后一项尤其值得注意——
模型自身的隐藏表示也被视作一种可被 tokenize 的模态，
这意味着上一轮推理的中间状态
可以作为下一轮推理的条件输入，
而无需写回自然语言。

Omni 的方法学贡献是 \textbf{Context Unrolling}。
传统多模态模型在产生预测前
通常只做一次前向传播，
跨模态推理被压缩进单层注意力的隐式过程；
Context Unrolling 则强制模型在产生预测前
\textbf{显式地、多步地跨多个模态表征做推理}，
每一步推理都会把当前结论以新 token 的形式追加到上下文。
这种「思考链发生在 token 层」的设计
让多模态推理的中间过程可观察、可复用、可被外部工具介入。

把 LongCat-Next 与 Omni 放到一起，
2026 年的统一 tokenizer 路线已经分化为两条主线。
LongCat-Next 强调
\textbf{tokenizer 本身的统一}——
让所有模态共享同一个离散字母表；
Omni 强调
\textbf{推理过程的统一}——
让所有模态共享同一个 unrolled 推理协议。
这两条路线并不互斥，
预计未来 12--18 个月会出现把两者合并的工作。

\subsection{VQ-VAE 的训练细节}

理解视觉分词需要了解 VQ-VAE 的核心机制。VQ-VAE（Vector Quantized Variational Autoencoder）的编码器将图像压缩为一组连续的潜在向量，然后通过\textbf{最近邻查找}将每个连续向量映射到码本（codebook）中最接近的离散码向量。这个量化步骤是视觉分词的关键，它将连续的图像表示转换为离散的 token 序列。

码本大小是视觉分词的核心超参数。Emu3 使用了 32,768 个码向量，这意味着每个视觉 token 可以表示 $\log_2(32768) = 15$ bits 的信息。一张 $32 \times 32 = 1024$ 个视觉 token 的图像因此包含了约 $1024 \times 15 = 15{,}360$ bits $\approx$ 1.9 KB 的信息——这远少于原始图像（一张 256$\times$256 的 RGB 图像约为 192 KB），说明 VQ-VAE 实现了约 100 倍的有损压缩。

这个压缩率与 BPE 对自然语言的压缩率类似：
BPE 将原始 UTF-8 字节序列压缩为更短的 token 序列，
压缩率约为 3--4$\times$（对英语）。
VQ-VAE 将原始像素值压缩为离散 token，
压缩率约为 100$\times$。
更高的压缩率意味着更大的信息损失：
这就是为什么当前的视觉 token 在重建细节上仍不如原始图像，
而文本 token 可以无损还原原始文本。
提升视觉 token 的保真度（在保持合理序列长度的同时减少信息损失）
是多模态分词的核心技术挑战之一。

一种有前景的方案是\textbf{残差量化}（Residual Quantization, RQ）：
不是用单个码向量表示一个 patch，
而是用多个码向量的「残差叠加」来逐步逼近原始表示。
每多一层残差，重建精度就提升一点。
这种方案允许在「精度」和「token 数」之间灵活权衡：
需要高精度时使用更多的残差层（更多 token），
需要高效率时使用更少的残差层（更少 token）。
SoundStream 和 EnCodec 在音频分词中已经大量使用了这种策略，
视觉分词领域也开始跟进。

VQ-VAE 训练的一个著名技术难点是\textbf{码本坍缩}（codebook collapse）——训练过程中，码本中大部分的码向量从未被使用，只有少数码向量承担了所有的编码工作。这类似于 BPE 词表中高频 token 和低频 token 的不平衡问题。常用的缓解策略包括指数移动平均更新码向量、随机重置不活跃的码向量、以及使用更大的码本结合正则化。

码本坍缩与 BPE 的 glitch token 问题有深层的对应关系：
两者都是由「频率不平衡」导致的：
高频的 token/码向量得到了充分训练，
低频的则沦为「死区」。
这再次说明了分词（无论是文本还是图像）中的一个核心矛盾：
\textbf{词表需要足够大来覆盖所有模式，
但词表中的每个条目又需要足够的训练信号来形成有用的表示}。
平衡这两个需求是分词器设计的永恒挑战。

对于文本分词，这个矛盾通过子词的组合性来缓解：
即使某个完整词不在词表中，
它的子词组合仍然可以提供有意义的表示。
但对于视觉分词，码向量之间缺乏类似的组合关系：
一个未训练的码向量就是一个「黑洞」，
与之关联的图像 patch 会产生低质量的重建。
如何让视觉分词也具备类似文本子词的「组合性」，
是多模态分词领域的一个重要开放问题。

一个可能的方向是\textbf{层次化码本}：
使用粗粒度的码向量捕获全局特征（如颜色、形状），
细粒度的码向量捕获局部细节（如纹理、边缘）。
粗粒度码向量的数量少，
可以获得充分的训练；
细粒度码向量的数量多，
但通过残差结构与粗粒度码向量共享信息。
这种层次化设计在概念上
类似于 BPE 的子词-词-短语的层次结构，
可能是实现视觉分词「组合性」的关键。

从更高的抽象层次看，文本分词和视觉分词正在走向融合。
BPE 将文本压缩为离散 token，VQ-VAE 将图像压缩为离散 token：
两者在形式上已经是统一的。
未来的统一多模态模型可能不再区分「文本分词器」和「视觉分词器」，
而是使用一个\textbf{通用的离散化框架}：
将任何连续信号（文本的字节流、图像的像素矩阵、音频的波形）
压缩为离散的 token 序列，
然后交由同一个 Transformer 处理。
这个愿景已经在 Emu3 和 Chameleon 中初步实现，
但距离真正的「统一分词器」还有很长的路要走。

\begin{corebox}[多模态分词的统一愿景]
多模态分词的发展正在走向一个统一的架构：
\textbf{所有模态共享同一个 token 空间}。
文本 token 和图像 token 混合在同一个序列中，
模型在预训练时同时学习所有模态的模式。
LLaMA 4 已经在词表中预留了多模态 token 的位置
（图像 token ID 为 200092），
Gemini 使用了类似的设计。
这种统一表示的终极目标是：
模型不再需要知道输入的模态是什么，
因为一切都已经是 token。
\end{corebox}

\section{分词的未来：仍是开放问题}

分词是 LLM 工程中最古老的组件之一，但它远未被「解决」。以下几个方向代表了分词领域最活跃的研究前沿。

\subsection{自适应分词}

BPE 的一个根本局限是：词表在训练后就固定了，无法根据输入内容动态调整。但不同类型的文本对分词的需求截然不同：处理数学论文时，数字和公式符号需要细粒度的分词；处理日常对话时，常见短语可以使用粗粒度的分词。这种「一刀切」的方式不是最优的。

考虑一个具体的场景：一段技术文档中混合了英文解释和 Python 代码片段。英文部分使用子词级别的分词效率最高（如 “tokenization” $\to$ “token” + “ization”），但代码部分可能需要字符级别的精细分词（如缩进空格需要逐个处理）。当前的 BPE 对两部分使用相同的分词策略，这意味着至少有一部分的分词不是最优的。

BLT 的动态 patch 分割是自适应分词的一种形式，它根据内容复杂度动态调整粒度。未来的分词器可能更进一步：在同一段文本中，对不同部分使用不同的粒度：代码块使用字符级别的精细分词，自然语言段落使用子词级别的高效分词，数字使用逐位分词。这种「混合粒度分词」在工程上的挑战是如何让 Transformer 高效地处理粒度不一致的 token 序列~\cite{pagnoni2024blt}。

一种可能的实现方案是\textbf{分层注意力}（hierarchical attention）：
模型首先在粗粒度的 token 层面做全局注意力（捕获长距离依赖），
然后在细粒度的子 token 层面做局部注意力（处理需要精细分析的片段）。
这种设计与人类阅读的模式一致：
我们先快速扫描全文获取大意，
然后对感兴趣的段落仔细阅读。

自适应分词的另一个动机来自\textbf{推理效率}。
在推理服务中，大部分请求是简单的日常对话：
这些请求用粗粒度的分词就足够了。
只有少部分请求涉及复杂的代码、数学或罕见语言：
需要细粒度的分词。
如果分词器能根据输入内容自动调整粒度，
就可以对简单请求使用更少的 token（降低计算成本），
对复杂请求使用更多的 token（保证质量）。
这种「按需分配」的策略与 MoE 架构的「按需激活」理念异曲同工。

\subsection{分词器的持续更新问题}

一个在实践中越来越突出的问题是：\textbf{分词器的词表在训练后就固定了，但语言是持续演变的}。

2024--2025 年涌现了大量新术语：“GRPO”、“MLA”、“vLLM”、“DeepSeek”，这些词在 2022 年的训练语料中不存在。如果分词器是 2022 年训练的，这些新词会被拆分为多个子词 token，降低效率。更根本的问题是：当模型通过 continual pre-training 或 RAG 获取新知识时，分词器的固定词表可能成为瓶颈。

一种解决方案是\textbf{词表扩展}，在不重新训练分词器的情况下，向词表中添加新的 token。这在技术上是可行的（只需扩展嵌入矩阵并初始化新 token 的向量），但引入了一个新问题：新 token 的嵌入需要通过额外的训练来学习，而这些新 token 在已有训练数据中不存在。实际上，LLaMA 3 从 LLaMA 2 的 32K 词表扩展到 128K 词表时，进行了完整的重新预训练——这是目前最可靠但也最昂贵的方案。

另一种思路是设计\textbf{模块化分词器}，将词表分为「核心词表」（固定不变的基础子词）和「扩展词表」（可以动态更新的高频模式）。核心词表保证了稳定性，扩展词表提供了灵活性。这个方向目前还在研究阶段，尚未被主流模型采用。

更激进的方案是\textbf{持续学习分词器}，即分词器在模型的整个生命周期中持续更新。当模型通过 continual pre-training 学习新知识时，分词器同步更新词表，将新领域的高频模式纳入词表。这种方案的技术挑战在于：词表的变化会导致已有 token 的嵌入失效，需要一种「增量更新」机制，在添加新 token 的同时保持旧 token 嵌入的稳定性。

一个折中的实现是 LLaMA 3 的方案：在版本迭代时（从 LLaMA 2 到 LLaMA 3）大幅扩展词表（32K $\to$ 128K），但在版本内部保持词表固定。这意味着词表的更新频率与模型的主要版本发布同步，通常是每 6--12 个月一次。对于大多数实际应用，这个更新频率是足够的。

但在快速变化的领域（如 AI 研究本身），
词汇的更新速度可能超过模型的迭代频率。
2024--2025 年涌现的大量新术语
（“GRPO”、“MLA”、“Sparse Attention”、“BLT”、“Engram”）
在当前模型的分词器中都不是独立 token：
它们会被拆分为多个子词。
这不影响模型对这些术语的理解
（模型可以从子词组合中推断含义），
但确实降低了效率。
一个有趣的研究方向是：
在保持模型其他参数不变的情况下，
\textbf{仅更新分词器和嵌入层}能带来多大的性能改善？
初步的实验表明，
这种「轻量级词表更新」可以以极低的成本
（通常只需要几千步的额外训练）
改善模型在新领域术语上的效率：
但对模型的通用能力影响可以忽略不计。
这种方法的实际价值在于：
它允许模型在部署后「适应」新的领域，
而无需进行完整的重新训练：
这对于需要持续跟踪前沿发展的应用场景
（如科研助手、技术文档检索）尤其有吸引力。

\subsubsection{字节级蒸馏：跨 tokenizer 的知识迁移}

分词器一旦被 BPE 锁死，
另一个长期被回避的问题是
\textbf{跨 tokenizer 的知识蒸馏}。
当教师模型（如 GPT-4 风格的 200K 词表）
和学生模型（如 LLaMA 系列的 128K 词表）
使用不同的分词器时，
两者的输出 logits 不在同一个概率单形上，
无法直接做 KL 散度对齐。
过去的解决方案大多是
在 token 序列层面做启发式对齐
（动态规划匹配公共子串、
近似投影到学生词表等），
工程复杂且容易在边界 token 上失真。

2026 年 4 月 MediaTek Research 与 UCL 提出的
\textbf{Byte-Level Distillation}（BLD）~\cite{liu2026bld}
给出了一个干净得多的方案：
用\textbf{字节级接口}作为公共底层。
具体做法是
把教师模型在每一步输出的
token 概率分布转换为字节级概率，
然后给学生模型加一个轻量的字节级 decoder head，
让学生在字节空间复现教师的字节级条件分布。
教师的 token 化方案、词表大小、特殊 token 集合
对学生完全透明——
学生只需要看见「下一个字节的概率」。

BLD 在 1B--8B 范围的实验中
匹配甚至超过了
此前最复杂的启发式 token 对齐基线，
而代码量、超参数数量都大幅下降。
更重要的是它给出了一个简洁的概念结论：
\textbf{字节级是跨 tokenizer 知识迁移的天然公约数}。
任何两个 UTF-8 兼容的语言模型，
不论词表如何，
都共享同一个字节空间；
只要愿意把对齐目标下沉到这一层，
跨 tokenizer 的训练信号就总能流通。
对开源社区而言，
这意味着不再被「教师与学生必须共用 tokenizer」
这一隐性约束绑住手脚——
任何强模型都可以作为蒸馏来源。

\subsection{可学习的分词}

传统分词器是与模型\textbf{分离训练}的，先训练分词器，再用固定的分词器训练模型。但这种分离意味着分词器无法获得模型的反馈。一个理想的方案是\textbf{端到端地联合训练分词器和模型}，让模型告诉分词器「什么样的切分方式最有利于我学习」。

这个方向的技术挑战在于：分词是一个离散操作（选择在哪里切分），而神经网络的反向传播需要连续的梯度。一些研究使用 Gumbel-Softmax~\cite{jang2017gumbel} 等技巧将离散的分词决策软化为连续的概率分布，从而实现端到端训练。但这些方法目前的训练效率还远不如传统的 BPE + 固定分词器流水线。

BLT 的动态 patch 分割可以被视为可学习分词的一种「弱形式」，虽然它不是端到端训练的（patch 边界由独立的熵模型决定），但它确实实现了「根据内容动态调整分词粒度」的目标。如果未来的研究能够将 BLT 的 patch 分割机制与 Transformer 的反向传播打通，让模型通过梯度信号告诉分词器「在这里切分更有利于我学习」——那将是分词领域的一个重大突破。

另一个可学习分词的思路是\textbf{hierarchical tokenization}，模型首先在字符或字节级别处理输入，然后通过一个可学习的「分组」模块将相邻的字符合并为更高层的「token」。这种分组是 soft 的（通过注意力权重实现），因此可以端到端训练。Perceiver 和 Perceiver IO 架构已经在这个方向上做了早期探索，用一组学习到的「查询」（query）从任意长度的输入中提取固定数量的「token」。将这种思路应用到语言模型中是一个活跃的研究方向。

\subsection{分词与模型坍缩}

2025 年出现了一个新的研究方向：分词器在合成数据（synthetic data）训练中的角色。当模型在自身生成的文本上训练时（所谓的「模型坍缩」），分词器的行为可能会放大这个问题。

原因是：模型生成的文本倾向于使用分词器中的高频 token（因为这些 token 的概率估计更准确），而回避低频 token。当这些生成的文本被用于下一轮训练时，低频 token 的训练信号进一步减少，形成一个正反馈循环，最终导致低频 token 的嵌入退化。这个问题在多语言场景中尤为严重：如果合成数据以英语为主，非英语 token 的嵌入质量会在迭代训练中逐渐下降。

识别和量化这种「分词器驱动的模型坍缩」是当前的活跃研究领域。一种可能的缓解策略是在合成数据生成中\textbf{刻意增加低频 token 的采样概率}，类似于人类在学习外语时刻意练习生词。

\subsection{分词与 MoE 架构的交互}

2025--2026 年 MoE 架构的广泛采用（DeepSeek-V3、Llama 4、Qwen3）引入了分词与架构之间的新交互效应。在 MoE 模型中，不同的专家（expert）被路由到不同的 token，路由决策通常基于 token 的隐层表示。

一个值得关注的现象是：\textbf{分词粒度会影响专家路由的行为}。如果一个中文四字词被编码为单个 token，它会被路由到一个专家处理。如果同一个词被拆分为 4 个 token，它们可能被路由到不同的专家，这意味着关于这个词的计算被分散到了多个专家之间，每个专家只看到了部分信息。DeepSeek-V3 的技术报告~\cite{deepseekv3} 提到了这个问题，并通过精心设计的路由策略来缓解，但更根本的解决方案是从分词层面确保语义完整的单元不被拆分。

\subsection{分词的标准化与互操作性}

当前 LLM 生态系统面临的一个实际问题是：\textbf{每个模型使用不同的分词器，导致 token 层面的工具和基础设施无法互通}。

一个提示模板在 ChatML 格式（Qwen）和 LLaMA 格式下产生不同的 token 序列。一段经过精心优化的 few-shot prompt 在更换模型后可能因为分词差异而失效。这种碎片化给应用开发者带来了显著的额外负担。

Hugging Face 的 \texttt{tokenizers} 库通过提供统一的 API 部分缓解了这个问题，开发者可以用相同的接口调用不同模型的分词器。但在更深层次上，分词器的标准化可能需要行业共识，类似于 Unicode 标准统一了字符编码。MCP（Model Context Protocol）的成功提供了一个先例：当一个标准获得足够的行业支持时，碎片化问题就能迅速解决。

一个值得关注的趋势是：越来越多的模型开始收敛到相似的分词方案：128K--150K 词表，tiktoken 风格的 BPE，支持多语言和代码。LLaMA 3、Qwen3 和 DeepSeek-V3 的分词器在设计理念上已经非常接近。如果这种收敛趋势继续，分词器的碎片化问题可能在未来几年内自然消解，不是因为行业达成了正式标准，而是因为最佳实践的趋同。

回到本章开头的类比：分词就像给模型配一副「眼镜」。
不同的分词器就像不同度数的镜片：
它们决定了模型能看清什么、看不清什么。
一副为英语优化的「眼镜」可能让模型对中文「近视」；
一副过于粗糙的「眼镜」可能让模型对数字「散光」。
选择正确的「眼镜」，这就是分词器设计的核心任务。

\subsection{分词在 Agent 场景中的新角色}

2025--2026 年 Agentic AI 的兴起给分词带来了新的考量。在 Agent 场景中，模型不仅处理自然语言，还需要与结构化的工具输入/输出交互，如 JSON 格式的函数调用、API 返回的结构化数据、文件系统路径、URL 等。

\textbf{结构化数据的分词效率}。一个典型的 JSON 工具调用可能如下：

\begin{lstlisting}[language={}]
{"function": "search", "arguments": {"query": "BPE tokenization"}}
\end{lstlisting}

这段文本中，大量的字符是 JSON 语法结构（花括号、引号、冒号），它们不携带任务相关的信息，但消耗了宝贵的 token 配额。一个为 Agent 场景优化的分词器可能将常见的 JSON 模式（如 \texttt{\{"function":}）合并为单个 token，从而在相同的上下文窗口内容纳更多有效信息。

\textbf{工具名称的 token 化}。Agent 频繁使用的工具名称（如 \texttt{search\_web}、\texttt{read\_file}、\texttt{execute\_code}）理想情况下应该是独立的 token，因为它们在 Agent 的决策链中是不可分割的语义单元。如果 \texttt{search\_web} 被拆分为 \texttt{search} + \texttt{\_} + \texttt{web}，模型需要额外的注意力计算来「重新组装」这个工具名称。Qwen3~\cite{qwen2025qwen3} 在词表中显式包含了常见的工具调用 token，使得模型可以更高效地处理 Agent 工作流。

\textbf{长上下文中的 token 预算管理}。在 Agent 场景中，上下文通常包含大量的历史对话、工具调用结果和系统提示，这些内容可能消耗上下文窗口的大部分配额。分词效率在这里变得尤为关键：一个 fertility 为 2.5 的分词器在 128K 上下文窗口中只能容纳约 51K 的有效「词」，而 fertility 为 1.2 的分词器可以容纳约 107K 的有效「词」，后者的有效容量是前者的 2 倍。

\begin{warnbox}[分词：简单表面下的深层影响]
本章的核心信息可以浓缩为一句话：\textbf{分词决定了模型「看到」的世界}。
从中文 fertility 到数学能力，从代码生成到多模态统一，
分词器的设计选择在模型的整个生命周期中持续产生影响：
而且这些影响往往是不可逆的：一旦分词器固定，
模型就被锁定在了这个特定的「感知方式」中。
这就是为什么分词值得一整章的深入讨论。

在 2026 年的技术栈中，分词仍然是「看不见但无处不在」的基础设施。
每一次 API 调用，每一轮模型训练，每一个 Agent 的工具调用：
都在分词器的定义下进行。
理解分词，就是理解 LLM 「感知世界」的方式。
这个理解对于做出正确的模型选择、优化推理成本、
以及诊断模型行为中的异常至关重要。

下一章（第~\ref{ch:data}章），
我们将讨论另一个「看不见但决定性」的因素：
训练数据的采集、清洗和工程。
如果分词是模型的「眼镜」，
那么数据就是模型的「教科书」。
两者共同决定了模型最终能学到什么。
\end{warnbox}


%% ============================================================
%% NEW REFERENCES (to be added to refs.bib):
%% gage1994bpe = Gage, Philip. A New Algorithm for Data Compression. C Users Journal, 12(2):23--38, 1994.
%% pennington2014glove = Pennington, Jeffrey and Socher, Richard and Manning, Christopher D. GloVe: Global Vectors for Word Representation. EMNLP, 2014.
%% schuster2012wordpiece = Schuster, Mike and Nakajima, Kaisuke. Japanese and Korean voice search. ICASSP, 2012.
%% rumbelow2023glitch = Rumbelow, Jessica and watMEG, Matthew. SolidGoldMagikarp (plus, prompt generation). Lesswrong, 2023.
%% mcleish2024abacus = McLeish, Sean et al. Transformers Can Do Arithmetic with the Right Embeddings. NeurIPS, 2024.
%% kreitner2024xval = Kreitner, Linus et al. Efficient numeracy in language models through single-token number embeddings. arXiv:2510.06824, 2024.
%% yu2023megabyte = Yu, Lili et al. MEGABYTE: Predicting Million-Byte Sequences with Multiscale Transformers. NeurIPS, 2023.
%% pagnoni2024blt = Pagnoni, Artidoro et al. Byte Latent Transformer: Patches Scale Better Than Tokens. arXiv:2412.09871, 2024.
%% wang2024emu3 = Wang, Xinlong et al. Emu3: Next-Token Prediction is All You Need. Nature, 2025.
%% ma2025unitok = Ma, Chuofan et al. UniTok: a Unified Tokenizer for Visual Generation and Understanding. NeurIPS, 2025.
%% xue2022byt5 = Xue, Linting et al. ByT5: Towards a Token-Free Future with Pre-trained Byte-to-Byte Models. TACL, 2022.
%% goldman2024tokenizer = Goldman, Omer et al. Unpacking Tokenization: Evaluating Text Compression and its Correlation with Model Performance. ACL Findings, 2024.
%% press2017tying = Press, Ofir and Wolf, Lior. Using the Output Embedding to Improve Language Models. EACL, 2017.
%% rust2021good = Rust, Phillip et al. How Good is Your Tokenizer? On the Monolingual Performance of Multilingual Language Models. ACL, 2021.
%% jang2017gumbel = Jang, Eric and Gu, Shixiang and Poole, Ben. Categorical Reparameterization with Gumbel-Softmax. ICLR, 2017.
%% cui2024llama = Cui, Yiming et al. "Efficient and Effective Text Encoding for Chinese LLaMA and Alpaca." arXiv:2304.08177, 2023.
%% deng2024multilingual = Deng, Yue et al. "Multilingual Jailbreak Challenges in Large Language Models." arXiv:2310.06474, 2024.
%% provilkov2020bpe = Provilkov, Ivan et al. "BPE-Dropout: Simple and Effective Subword Regularization." ACL, 2020.
%% zheng2024judging = Zheng, Lianmin et al. "Judging LLM-as-a-Judge with MT-Bench and Chatbot Arena." NeurIPS 2023.
%% ============================================================
